\section{The twisted places}\label{sec:twisted}

At the places where the local type of the lifts is twisted, the lattice of integral lifted forms differs from the lattice that the period pairing sees, and a congruence module measures the difference. This section defines the congruence defects $\eta_j(w)$, states the two bounds that control them, and proves the parts of those bounds that reduce to commutative algebra and Galois cohomology. Theorem~\ref{r5:thm:DIVinN} combines the two bounds. Its conclusion enters the heights of the period maps in \eqref{A4} and the frame parameters in \eqref{A3}. It does not enter \eqref{A5}, which Corollary~\ref{r5:cor:A5} proves without it.

The first bound, Hypothesis~\ref{w4:thm:vwin}, bounds the lengths of Tate--Shafarevich groups of the host field on windows of scales. Subsection~\ref{tw:sec:vandiver} shows that it implies a bound in the direction of Vandiver's conjecture beyond the bounds published by Soul\'e~\cite{soule99} and by de~la~Torre~Dur\'an~\cite{delatorre}, and traces the argument for it to an automorphy statement for residually reducible Galois representations that is not a published theorem.

\subsection{Congruence defects}\label{tw:sec:eta}

For each finite place $w$ of $L$ the lifts have a local type $\mu_w$, which is either the standard type $\mu_0$ or a twisted type $\mu^{\rm tw}_w$, and $\mu_{\bar w}=\mu_w$. We call $w$ \emph{twisted} when $\mu_w\ne\mu_0$. At a twisted place the integral structure of the twisted lifts at level $j$ differs from the integral structure of the eigen-components that the period pairing sees, and the difference is measured by a congruence module.

\begin{definition}\label{r5:def:eta}
Let $w$ be a twisted place, $l=p_w$, and $1\le j\le m$. Let $\mathcal O$ be the ring of integers of the completion of $L$ at $w$, let $\Lambda$ be the $\mathcal O$-lattice spanned by the integral forms of level $j$ and local type $\mu^{\rm tw}_w$, and let $\mathbb T$ be the $\mathcal O$-algebra of Hecke operators acting on $\Lambda$. For a maximal ideal $\mathfrak m$ of $\mathbb T$ split $\Lambda_{\mathfrak m}\otimes\mathbb Q$ into the sum $V_{\rm good}\oplus V_{\rm bad}$ of the eigen-components whose eigensystems occur in the lifts of level $j$ and of the others. Put
\[
\eta_j(w)=\sum_{\mathfrak m}\frac{\operatorname{length}_{\mathcal O}\Big(\big(V_{\rm good}\cap(\Lambda_{\mathfrak m}+V_{\rm bad})\big)\big/\big(\Lambda_{\mathfrak m}\cap V_{\rm good}\big)\Big)}{\operatorname{length}_{\mathcal O}(\mathcal O/l)} ,
\]
the normalized length of the congruence module of $V_{\rm good}$ in $\Lambda$. For a non-twisted place put $\eta_j(w)=0$.
\end{definition}

The module $(V_{\rm good}\cap(\Lambda+V_{\rm bad}))/(\Lambda\cap V_{\rm good})$ is the projection of $\Lambda$ to $V_{\rm good}$ modulo the forms of $\Lambda$ that lie in $V_{\rm good}$. It vanishes when no good eigensystem is congruent to a bad one modulo $\mathfrak m$. For $1\le j\le m$ we write $d^0_j$ for the rank of the space of lifted forms of level $j$ and standard type $\mu_0$, and $N_c=\lvert T_2\rvert\,k\,b\,x^{2\varepsilon}/n^2$, so that $\mathcal B=N_Ta(1+m^2/k)+m(B_0+N_c)$.

\subsection{Tate--Shafarevich groups}\label{tw:sec:sha}

Let $F$ be a number field, $\overline F$ an algebraic closure, $G_F=\Gal(\overline F/F)$ with the Krull topology, and $l$ a prime. Let $\chi\colon G_F\to\Zl^\times$ be the $l$-adic cyclotomic character, so that $\sigma(\zeta)=\zeta^{\chi(\sigma)}$ for every $l$-power root of unity $\zeta$. For an integer $e$ let $M_e$ be the discrete group $\Ql/\Zl$ on which $\sigma$ acts by multiplication by $\chi(\sigma)^{-e}$. For a finite place $v$ of $F$ fix a place of $\overline F$ above it, with decomposition group $D_v$ and inertia group $I_v$. Cohomology is continuous cohomology with discrete coefficients.

\begin{definition}\label{w4:def:sha}
$\Sha_l(F,e)$ is the subgroup of $H^1(G_F,M_e)$ of classes $c$ with $c|_{D_v}=0$ for every $v\mid l$ and $c|_{I_v}=0$ for every finite $v$. For a finite free $\Zl$-algebra $O$ put $M_{e,O}=O\otimes_{\Zl}M_e$, with $G_F$ acting on the second factor, and define $\Sha_l(F,e;O)\subset H^1(G_F,M_{e,O})$ by the same conditions. For a finite abelian $l$-group $X$ write $\lambda(X)=\log|X|=\len_{\Zl}(X)\log l$.
\end{definition}

The conditions do not depend on the chosen places above $v$, because restriction to conjugate subgroups differs by an automorphism of cohomology. There is no condition at archimedean places. For odd $l$ none is needed, because $H^1$ of a group of order $2$ with coefficients in an $l$-primary group vanishes.

\begin{lemma}\label{w4:lem:unram}
Let $v\nmid l$ be a finite place of $F$ and $e\ge1$. An element of $H^1(D_v,M_e)$ that vanishes on $I_v$ is zero.
\end{lemma}

\begin{proof}
The group $I_v$ acts trivially on $M_e$, because $\chi$ is unramified at $v$. Inflation--restriction gives an exact sequence $0\to H^1(D_v/I_v,M_e)\to H^1(D_v,M_e)\to H^1(I_v,M_e)$, so a class that vanishes on $I_v$ comes from $H^1(D_v/I_v,M_e)$. The group $D_v/I_v$ is procyclic, topologically generated by a Frobenius element $\varphi$, and $M_e$ is a discrete torsion module, so $H^1(D_v/I_v,M_e)\cong M_e/(\varphi-1)M_e$. The element $\varphi$ acts on $M_e$ as multiplication by $a=Nv^{-e}$. Since $Nv>1$ and $e\ge1$, $a-1$ is a nonzero element of $\Ql$, say $a-1=l^ku$ with $u\in\Zl^\times$. Multiplication by $l^ku$ is surjective on the divisible group $\Ql/\Zl$, so $M_e/(\varphi-1)M_e=0$.
\end{proof}

\begin{lemma}\label{w4:lem:indepS}
Let $e\ge1$ and let $S$ be a finite set of finite places of $F$ containing the places above $l$. Let $F_S\subset\overline F$ be the maximal extension of $F$ unramified outside $S$ and the archimedean places, and $G_{F,S}=\Gal(F_S/F)$. Inflation identifies $\Sha_l(F,e)$ with
\[
\ker\Bigl(H^1(G_{F,S},M_e)\longrightarrow\prod_{v\in S}H^1(D_v,M_e)\Bigr).
\]
In particular $\Sha_l(F,e)$ does not depend on $S$.
\end{lemma}

\begin{proof}
The module $M_e$ is unramified outside $l$, so $H=\Gal(\overline F/F_S)$ acts trivially on it. Inflation--restriction gives $0\to H^1(G_{F,S},M_e)\to H^1(G_F,M_e)\to\Hom(H,M_e)$. The group $H$ is the closed normal subgroup of $G_F$ generated by the inertia groups at places outside $S$. A class $c\in H^1(G_F,M_e)$ that vanishes on every $I_v$ restricts to a homomorphism on $H$ that is trivial on these generators and is invariant under conjugation by $G_F$ (as an element of $\Hom(H,M_e)^{G_F}$), hence trivial on $H$. So $c$ is inflated from $G_{F,S}$. Conversely an inflated class vanishes on $I_v$ for $v\notin S$. At $v\notin S$ an unramified class is zero by Lemma~\ref{w4:lem:unram}. At $v\in S$ with $v\nmid l$, vanishing on $D_v$ implies vanishing on $I_v$, and vanishing on $I_v$ implies vanishing on $D_v$ by Lemma~\ref{w4:lem:unram}. At $v\mid l$ both definitions impose vanishing on $D_v$, which implies vanishing on $I_v$. This proves the identification, and it shows that the group does not depend on $S$.
\end{proof}

\begin{lemma}\label{w4:lem:coef}
Let $O$ be a finite free $\Zl$-algebra which is a discrete valuation ring with ramification index $e(O)$ over $\Zl$. Then $\Sha_l(F,e;O)\cong O\otimes_{\Zl}\Sha_l(F,e)$ as $O$-modules, and if $\Sha_l(F,e)$ is finite,
\[
\len_O\Sha_l(F,e;O)=e(O)\,\len_{\Zl}\Sha_l(F,e),\qquad \len_O(O/lO)=e(O).
\]
\end{lemma}

\begin{proof}
Choose a $\Zl$-basis $o_1,\dots,o_r$ of $O$. Then $M_{e,O}=\bigoplus_i o_i\otimes M_e$ as $G_F$-modules, so $H^1(G_F,M_{e,O})=\bigoplus_iH^1(G_F,M_e)$, compatibly with restriction to every closed subgroup. Hence $\Sha_l(F,e;O)=\bigoplus_io_i\otimes\Sha_l(F,e)=O\otimes_{\Zl}\Sha_l(F,e)$, and the $O$-module structures agree because $O$ acts on the first tensor factor in both. The functor $X\mapsto O\otimes_{\Zl}X$ is exact, and it sends $\Z/l$ to $O/lO$, whose $O$-length is $e(O)$. By additivity of length along a composition series of $X$, $\len_O(O\otimes X)=e(O)\len_{\Zl}X$ for every $\Zl$-module $X$ of finite length.
\end{proof}

\begin{lemma}\label{w4:lem:pm}
Let $l$ be odd, $E$ a CM field, $c\in\Gal(E/\Q)$ complex conjugation, $e\ge1$. The natural action of $c$ on $H^1(G_E,M_e)$ preserves $\Sha_l(E,e)$ and squares to the identity. If $\Sha_l(E,e)$ has finite length then $\len\Sha_l(E,e)=\len\Sha_l(E,e)^{c=(-1)^e}+\len\Sha_l(E,e)^{c=(-1)^{e+1}}$.
\end{lemma}

\begin{proof}
The natural action of $\sigma\in G_\Q$ on $H^1(G_E,M_e)$ is $(\sigma\cdot f)(g)=\sigma f(\sigma^{-1}g\sigma)$ on cocycles. It factors through $\Gal(E/\Q)$, permutes the decomposition and inertia groups, and so preserves $\Sha_l(E,e)$. The action of $c$ is an involution. Since $2$ is a unit in $\Zl$, $x=\frac12(x+cx)+\frac12(x-cx)$ splits $\Sha_l(E,e)$ as the direct sum of the $c$-eigenspaces for $\pm1$, and length is additive.
\end{proof}

\begin{proposition}[Restriction to abelian fields]\label{w4:prop:res}
Let $N\ge1$, let $E\subset\Q(\mu_N)$ be a subfield, let $l\nmid N$ be an odd prime and let $e\ge1$ with $(l-1)\nmid e$. Restriction $H^1(G_\Q,M_e)\to H^1(G_E,M_e)$ maps $\Sha_l(\Q,e)$ injectively into $\Sha_l(E,e)$, onto classes fixed by $\Gal(E/\Q)$. If $E$ is a CM field and $e$ is even, the image lies in $\Sha_l(E,e)^{c=(-1)^e}$. In particular $\len\Sha_l(\Q,e)\le\len\Sha_l(E,e)$.
\end{proposition}

\begin{proof}
Restriction of a class that vanishes on $D_v$ (respectively $I_v$) vanishes on $D_w=D_v\cap G_E$ (respectively $I_w$) for every place $w$ of $E$ above $v$, so the image lies in $\Sha_l(E,e)$. The kernel of restriction is $H^1(\Gal(E/\Q),M_e^{G_E})$. Since $l\nmid N$, the fields $\Q(\mu_N)$ and $\Q(\mu_{l^\infty})$ are linearly disjoint over $\Q$, because a subfield of both is unramified at every prime. Hence $\chi(G_E)\supseteq\chi(G_{\Q(\mu_N)})=\Zl^\times$. An element $x\in M_e[l]$ fixed by $G_E$ satisfies $u^{-e}x=x$ for all $u\in\Zl^\times$, and since $(l-1)\nmid e$ there is $u$ with $u^{-e}\not\equiv1\pmod l$, so $x=0$. A nonzero $G_E$-fixed element of $M_e$ of order $l^k$ would give the nonzero fixed element $l^{k-1}x$ of $M_e[l]$. So $M_e^{G_E}=0$ and restriction is injective. A restricted class $f=c|_{G_E}$ satisfies $(\sigma\cdot f)(g)=\sigma c(\sigma^{-1}g\sigma)=c(g)$ modulo coboundaries, by the cocycle relation for $c$, so it is fixed by $\Gal(E/\Q)$. For $E$ CM and $e$ even, $c$ acts trivially on the image, and $1=(-1)^e$.
\end{proof}

\begin{proposition}[Class groups]\label{w4:prop:classgroup}
Let $l$ be an odd prime, $e\ge1$ with $(l-1)\nmid e$, $F_l=\Q(\mu_l)$, $\Delta=\Gal(F_l/\Q)$, $\omega\colon\Delta\to\F_l^\times$ the Teichm\"uller character, and $A=\Cl(F_l)\otimes\Zl$. For a character $\theta$ of $\Delta$ let $A^{(\theta)}=\{a\in A:\delta a=\theta(\delta)a\text{ for all }\delta\}$. Then
\[
\dim_{\F_l}\Sha_l(\Q,e)[l]=\dim_{\F_l}\bigl(A/lA\bigr)^{(\omega^{-e})}.
\]
In particular $\Sha_l(\Q,e)=0$ if and only if $A^{(\omega^{-e})}=0$.
\end{proposition}

\begin{proof}
\emph{Step 1.} $H^0(G_\Q,M_e)=0$ and $H^0(D_l,M_e)=0$. The group $D_l$ maps onto $\Delta$ (the prime $l$ is totally ramified in $F_l$) and acts on $M_e[l]$ through $\omega^{-e}$, which is nontrivial because $(l-1)\nmid e$. So $M_e[l]^{D_l}=0$, and as in the proof of Proposition~\ref{w4:prop:res} this gives $M_e^{D_l}=0$, and a fortiori $M_e^{G_\Q}=0$.

\emph{Step 2.} From $0\to M_e[l]\to M_e\xrightarrow{l}M_e\to0$ and Step 1, the map $H^1(G_\Q,M_e[l])\to H^1(G_\Q,M_e)[l]$ is an isomorphism. For $v\ne l$ the group $I_v$ acts trivially on $M_e$, so the kernel of $H^1(I_v,M_e[l])\to H^1(I_v,M_e)$ is $M_e/lM_e=0$. At $l$ the kernel of $H^1(D_l,M_e[l])\to H^1(D_l,M_e)$ is $H^0(D_l,M_e)/l=0$. Hence $\Sha_l(\Q,e)[l]$ is identified with the group $\mathcal S$ of $c\in H^1(G_\Q,\F_l(\omega^{-e}))$ that vanish on $I_v$ for $v\ne l$ and on $D_l$.

\emph{Step 3.} Since $|\Delta|=l-1$ is prime to $l$, restriction gives an isomorphism $H^1(G_\Q,\F_l(\omega^{-e}))\cong H^1(G_{F_l},\F_l(\omega^{-e}))^\Delta=\Hom(G_{F_l},\F_l)^{\Delta}$, where $G_{F_l}$ acts trivially on the coefficients and $\Delta$-invariance of $f$ means $f(\tilde\delta g\tilde\delta^{-1})=\omega^{-e}(\delta)f(g)$. For $v\ne l$ the extension $F_l/\Q$ is unramified at $v$, so $I_v\subset G_{F_l}$, and $c$ vanishes on $I_v$ iff $f$ does. Restriction from $D_l$ to $D_l\cap G_{F_l}=D_\lambda$ is injective on $H^1(-,\F_l(\omega^{-e}))$ (index prime to $l$), so $c|_{D_l}=0$ iff $f|_{D_\lambda}=0$, where $\lambda=(1-\zeta_l)$ is the prime above $l$.

\emph{Step 4.} The field $F_l$ is totally imaginary. By class field theory a homomorphism $G_{F_l}\to\F_l$ that vanishes on every inertia group factors through the Artin isomorphism $\Art\colon\Cl(F_l)\to\Gal(H/F_l)$ of the Hilbert class field $H$. Its restriction to $D_\lambda$ is determined by its value on $\Art([\lambda])$, which is $0$ because $\lambda$ is principal. The Artin map satisfies $\tilde\delta\Art(a)\tilde\delta^{-1}=\Art(\delta a)$. So $\mathcal S\cong\Hom_\Delta(A/lA,\F_l(\omega^{-e}))$, homomorphisms with $f(\delta a)=\omega^{-e}(\delta)f(a)$. Since $|\Delta|$ is prime to $l$, $A/lA=\bigoplus_\theta(A/lA)^{(\theta)}$, and such an $f$ vanishes on $(A/lA)^{(\theta)}$ for $\theta\ne\omega^{-e}$. Hence $\mathcal S\cong\Hom_{\F_l}((A/lA)^{(\omega^{-e})},\F_l)$. Finally $\Sha_l(\Q,e)$ is an $l$-primary torsion group, so it vanishes iff its $l$-torsion does, and $(A/lA)^{(\theta)}=A^{(\theta)}/lA^{(\theta)}$ vanishes iff $A^{(\theta)}$ does.
\end{proof}

For even $e$ the character $\omega^{-e}$ is even. Vandiver's conjecture for $l$ is the vanishing of $A^{(\theta)}$ for all even $\theta$. By Proposition~\ref{w4:prop:classgroup} it is equivalent to $\Sha_l(\mathbb Q,e)=0$ for all even $e$ with $2\le e\le l-3$, since the trivial character contributes $0$.

For comparison we record the following published identification. Let $l$ be odd and $e\ge1$ with $(l-1)\nmid e$. By Lemma~\ref{w4:lem:indepS} with $S=\{l\}$ and Poitou--Tate duality~\cite[Theorem~I.4.10]{milneADT}, $\Sha_l(\mathbb Q,e)$ is Pontryagin dual to $H^2(\mathbb Z[1/l],\Zl(1+e))$, and this group is $K_{2e}(\mathbb Z)\otimes\Zl$ by the theorem of Rost and Voevodsky on the Quillen--Lichtenbaum conjecture for $\mathbb Z$ (see the survey of Weibel~\cite{weibelK}). Nothing below depends on this identification.

\subsection{The Selmer bound on windows}\label{tw:sec:vwin}

\begin{definition}[Window pairs]\label{w4:def:winpair}
Let $0<\varepsilon<\frac1{52}$ and $Y\ge2$. A \emph{window pair at scale $Y$} is a pair of primes $(q,d)$ with $q\equiv3\pmod8$, $\frac12Y^{\varepsilon/5}<\log q\le Y^{\varepsilon/5}$, $d\equiv7\pmod{16}$, $\frac12Y^{\varepsilon/5}\log^2Y<\log d\le Y^{\varepsilon/5}\log^2Y$, $q\mid d+1$, and $p\mid d+1$ for every odd prime $p\le Y^{\varepsilon/5}$.
\end{definition}

\begin{definition}[$\chi$-fields]\label{w4:def:chifield}
Let $\psi_0$ be a primitive Dirichlet character of conductor $N_0$, $e\ge0$, and $q,d$ odd primes. Put $N=\operatorname{lcm}(2^{e+2},N_0,q,d)$. The field $E(\psi_0,e,q,d)$ is the subfield of $\Q(\mu_N)$ fixed by the automorphisms $\zeta\mapsto\zeta^a$ with $a\equiv\pm1\pmod{2^{e+2}}$, $\psi_0(a)=1$, $\left(\frac aq\right)=\left(\frac ad\right)=1$. For $q\equiv3\pmod4$ it is a CM field.
\end{definition}

\begin{hypothesis}[Selmer bound on windows]\label{w4:thm:vwin}
For every primitive Dirichlet character $\psi_0$, every $0<\varepsilon<\frac1{52}$ and every function $e\colon[1,\infty)\to\mathbb N$ with $x^{2+8\varepsilon}\le2^{e(x)}\le2x^{2+8\varepsilon}$, there are $J\ge1$ and $X_0$ with the following property. For every $X\ge X_0$ and every window pair $(q,d)$ at scale $X^{1/J}$ there is $Y$ with $X^{1/J}\le Y\le X^{1/2}$ such that for every $x\in[Y,Y^2]$, with $t=\lfloor x^{401}\rfloor$ and $E=E(\psi_0,e(x),q,d)$, every $\Sha_l(E,e')$ with $0<e'<t$ is finite and
\[
\sum_{l}\ \sum_{1\le e'<t}\lambda\bigl(\Sha_l(E,e')\bigr)\le\exp\bigl(\tfrac12x^{\varepsilon}\bigr),
\]
the sum over all primes $l$ (including $l=2$).
\end{hypothesis}

Proposition~\ref{r6:R6:thm:selmer} deduces (D6) by rescaling the window.

For a CM field $E$, $t\ge1$ and $B\ge0$ write
\begin{itemize}[leftmargin=*,itemsep=1pt]
\item $(\mathrm P_2)(E,t,B)$: $\sum_{1\le e'<t}\len_{\mathbb Z_2}\Sha_2(E,e')\le B$, all finite;
\item $(\mathrm P_\pm)(E,t,B)$: for odd $l$, the $c=\pm(-1)^{e'}$ parts are finite and
\[
\sum_{l\ \mathrm{odd}}\log l\sum_{e'<t}\len\Sha_l(E,e')^{c=\pm(-1)^{e'}}\le B.
\]
\end{itemize}

\begin{theorem}[decomposition]\label{w4:thm:decomp}
Hypothesis~\ref{w4:thm:vwin} follows from the following two statements, with the notation of Hypothesis~\ref{w4:thm:vwin}.
\begin{enumerate}[leftmargin=*,label=(\Roman*)]
\item \emph{(Plus part and $l=2$ on windows.)} For every $\psi_0$, $\varepsilon$ and every pair of functions $e(x),t(x)$ with $2\le e(x)$, $2^{e(x)}\le x^4$, $2\le t(x)\le x^{401}$ for $x\ge2$, there are $J,C,P,X_1>0$ such that for $X\ge X_1$ and primes $q\equiv3\ (8)$, $d\equiv7\ (16)$ prime to $N_0$, with $q,d>\max(P,2X^{401}+1)$ and $\log(qd)\le X^{\varepsilon/(4J)}$, there is $Y$ with $X^{1/J}\le Y\le X^{1/2}$ such that for all $x\in[Y,Y^2]$ and $E=E(\psi_0,e(x),q,d)$: $(\mathrm P_2)(E,t(x),e^{x^{4\varepsilon/5}})$ and $(\mathrm P_+)(E,t(x),(\cond(E)t(x))^C+Ce^{x^{4\varepsilon/5}})$.
\item \emph{(Minus part.)} There is $C$ such that for all such $E$ of degree $n$ and conductor $f$, and $t\ge2$: $(\mathrm P_-)(E,t,\beta)$ with $\beta=3nt^2\log(tf)+C\sum_{l\mid o,\ l\ \mathrm{odd}}l\,n^2t^2\log^2(tlf)$, where $o$ is the order of $\psi_0$.
\end{enumerate}
\end{theorem}

\begin{proof}
Fix $\psi_0$, $\varepsilon$ and $e(x)$ as in Hypothesis~\ref{w4:thm:vwin}, and put $t(x)=\lfloor x^{401}\rfloor$. For $x\ge2$, $2^{e(x)}\le2x^{2+8\varepsilon}\le x^4$ and $2^{e(x)}\ge x^2\ge4$, so (I) applies; let $J_{\rm I},C,P,X_1$ be its constants and put $J=\max(J_{\rm I},1)$, and let $C_-$ be the constant of (II).

\emph{Sizes.} Let $(q,d)$ be a window pair at scale $Y_0=X^{1/J}$ and $x\ge Y_0$. Then $\log q+\log d\le Y_0^{\varepsilon/5}(1+\log^2Y_0)\le x^{\varepsilon/4}$ once $Y_0$ is large. Also $q,d\ge\exp(\frac12Y_0^{\varepsilon/5})>\max(P,2X^{401}+1)$ and $\log(qd)\le X^{\varepsilon/(4J)}$ for $X$ large, so the hypotheses of (I) hold. The field $E$ lies in $\Q(\mu_N)$ with $N=\operatorname{lcm}(2^{e(x)+2},N_0,q,d)$, so $f=\cond(E)\le N$ and $n=[E:\Q]\le\varphi(N)\le N$. Since $2^{e(x)}\le2x^{3}$,
\[
\log(ft)\le(e(x)+2)\log2+\log N_0+x^{\varepsilon/4}+401\log x\le2x^{\varepsilon/4}
\]
for $x$ large in terms of $\varepsilon$ and $N_0$.

\emph{Assembly.} Let $x\in[Y,Y^2]$ with $Y$ from (I). For odd $l$, Lemma~\ref{w4:lem:pm} splits each length into the two parts. Summing over $l$,
\[
\sum_l\sum_{e'<t}\lambda(\Sha_l(E,e'))\le e^{x^{4\varepsilon/5}}\log2+\beta+(ft)^C+Ce^{x^{4\varepsilon/5}}.
\]
Since $n\le f$, $\ell\le o$ for $\ell\mid o$, and $\log^2y\le y$ for $y\ge1$, each summand of $\beta$ is at most $o^2(ft)^3$, so $\beta\le(3+C_-o^3)(ft)^3\le(3+C_-o^3)e^{6x^{\varepsilon/4}}$; also $(ft)^C\le e^{2Cx^{\varepsilon/4}}$. Each of the four terms is at most $\frac14\exp(\frac12x^\varepsilon)$ once $x$ is large in terms of $\varepsilon$, $C$, $C_-$, $o$, because $\frac45\varepsilon<\varepsilon$ and $\frac14\varepsilon<\varepsilon$. Choose $X_0\ge X_1$ so large that $Y_0=X_0^{1/J}$ exceeds all thresholds above. Then for $X\ge X_0$ every $x\ge Y\ge X^{1/J}$ is large enough.
\end{proof}

\notgiven{Neither of the statements (I) and (II) is proved here.} The plus part of (I) splits, by the residual position of $l$ relative to the host, into five counts of the same type. The generic one is a bound for the groups $\Sha_l(\mathbb Q,e';\psi)$ attached to the Dirichlet characters $\psi$ of the host, and it follows from statement (CorGK) of Subsection~\ref{tw:sec:vandiver} for $K=\mathbb Q(\sqrt{-p_K})$. Cases (II) and (III) use $\mathrm{GSp}_4$. Case (V) uses a deformation problem over a frame of the host field. The part of (I) at $l=2$ and at the bad primes uses Poitou--Tate duality, local duality, a surjection onto $H^2(G_{L,S},\mathbb Z_l(e+1))$ from a module that maps to $\mathbb Z_l\otimes H_{2e}(\mathrm{GL}_{4e+4}(\mathcal O_L),\mathbb Z)$ with kernel killed by $l^{c(2e)}$, torsion bounds of the type of Soul\'e and of Dwyer and Friedlander, and Borel's computation of ranks. The bound is available only on windows of scales, and the constant $J=2^{R+2}(2C_K\cdot401/\varepsilon)^R$ of the formal proof, with $R=R(\psi_0)$, depends on $\psi_0$ and $\varepsilon$. Statement (II) is of the type of the Main Conjecture for abelian fields (Mazur and Wiles, Wiles, Rubin) together with the theorem of Ferrero and Washington, which bound the minus parts by generalized Bernoulli numbers, whose logarithms are polynomial in $n$, $t$ and $\log f$. The formal proof derives (II) from Stickelberger's theorem and a Kolyvagin system of Gauss sums.

\begin{proposition}[plus parts through characters]\label{w4:prop:pluschar}
Let $N\ge1$, let $H\subset(\Z/N)^\times$ be a subgroup with $-1\notin H$, and let $E\subset\Q(\mu_N)$ be the fixed field of $H$, so that $E$ is a CM field with $\Gal(E/\Q)\cong(\Z/N)^\times/H$. Let $l$ be an odd prime with $l\nmid[(\Z/N)^\times:H]$, let $e\ge1$, and let $O$ be a finite extension of $\Zl$ which is a discrete valuation ring and contains the values of all characters of $(\Z/N)^\times/H$. Let $X_e$ be the set of Dirichlet characters $\psi$ modulo $N$ with values in $O$ that are trivial on $H$ and satisfy $\psi(-1)=(-1)^e$. Then $\Sha_l(E,e)^{c=(-1)^e}$ is finite if and only if every $\Sha_l(\Q,e;\psi)$ with $\psi\in X_e$ is finite, and in that case
\[
\len_{\Zl}\Sha_l(E,e)^{c=(-1)^e}=\sum_{\psi\in X_e}\frac{\len_O\Sha_l(\Q,e;\psi)}{e(O)} .
\]
\end{proposition}

\begin{proof}
Put $\Gamma=\Gal(E/\Q)$, $M=M_{e,O}$ and $A=O[\Gamma]$, on which $G_\Q$ acts through $\Gamma$ by left multiplication. Since $M$ is a $G_\Q$-module, the induced module $\operatorname{Ind}_{G_E}^{G_\Q}M$ is isomorphic to $A\otimes_OM$ with the diagonal action, and Shapiro's lemma gives an isomorphism $H^1(G_E,M)\cong H^1(G_\Q,A\otimes_OM)$. The order of $\Gamma$ is a unit in $O$ and $O$ contains the values of the characters of $\Gamma$, so the elements $\varepsilon_\psi=\lvert\Gamma\rvert^{-1}\sum_{\gamma\in\Gamma}\psi(\gamma)^{-1}\gamma$ are orthogonal idempotents with $\sum_\psi\varepsilon_\psi=1$ and $\gamma\varepsilon_\psi=\psi(\gamma)\varepsilon_\psi$. Hence $A\otimes_OM=\bigoplus_\psi\varepsilon_\psi A\otimes_OM\cong\bigoplus_\psi M(\psi)$ as $G_\Q$-modules, the sum over all characters $\psi$ of $\Gamma$, where $M(\psi)$ is $M$ with the action twisted by $\psi$. This gives $H^1(G_E,M)\cong\bigoplus_\psi H^1(G_\Q,M(\psi))$.

Let $v$ be a finite prime and let $D$ be $D_v$ or $I_v$. By Mackey's formula the restriction of $\operatorname{Ind}_{G_E}^{G_\Q}M$ to $D$ is the sum over the double cosets $D\gamma G_E$ of the modules induced from $D\cap\gamma G_E\gamma^{-1}$, and Shapiro's lemma for each summand identifies the restriction map to $H^1(D,A\otimes_OM)$ with the product of the restriction maps of $H^1(G_E,M)$ to the decomposition groups, respectively the inertia groups, of the places of $E$ above $v$. So a class of $H^1(G_E,M)$ lies in $\Sha_l(E,e;O)$ if and only if its image in $H^1(G_\Q,A\otimes_OM)$ vanishes on $D_l$ and on every $I_v$. The decomposition $A\otimes_OM=\bigoplus_\psi M(\psi)$ is stable under $G_\Q$, hence under $D$, so these conditions hold if and only if they hold in each component. Therefore $\Sha_l(E,e;O)\cong\bigoplus_\psi\Sha(\psi)$, where $\Sha(\psi)\subset H^1(G_\Q,M(\psi))$ is the subgroup of classes vanishing on $D_l$ and on every $I_v$. Since $\psi\mapsto\psi^{-1}$ is a bijection of the characters of $\Gamma$, $\Sha(\psi)$ is the group $\Sha_l(\Q,e;\psi^{-1})$ of the text.

The natural action of $\gamma\in\Gamma$ on $H^1(G_E,M)$ of Lemma~\ref{w4:lem:pm} corresponds under Shapiro's isomorphism to the action of $A\otimes_OM$ induced by multiplication by $\gamma$ or by $\gamma^{-1}$ on $A$, according to the convention, which commutes with $G_\Q$ because $\Gamma$ is abelian. On $\varepsilon_\psi A$ it is multiplication by $\psi(\gamma)^{\pm1}$. Complex conjugation $c$ is the class of $-1$, and $\psi(-1)^{\pm1}=\psi(-1)$. Hence $\Sha_l(E,e;O)^{c=(-1)^e}\cong\bigoplus_{\psi(-1)=(-1)^e}\Sha(\psi)$. The set $X_e$ is stable under $\psi\mapsto\psi^{-1}$, so the right side is $\bigoplus_{\psi\in X_e}\Sha_l(\Q,e;\psi)$. Since $2$ is a unit in $O$, the eigenspace decomposition commutes with $O\otimes_{\Zl}-$, and Lemma~\ref{w4:lem:coef} gives $\len_O\Sha_l(E,e;O)^{c=(-1)^e}=e(O)\len_{\Zl}\Sha_l(E,e)^{c=(-1)^e}$. Length is additive on finite direct sums, which gives both assertions.
\end{proof}

\begin{proposition}[the generic plus part from (CorGK)]\label{w4:prop:genericplus}
Assume statement (CorGK) of Proposition~\ref{w4:prop:corToBound} for an imaginary quadratic field $K$, with constants $c_5,A$. Let $N\ge16$ be an integer, $e\ge1$ and $t\ge e+1$ integers, $r_1\ne r_2$ primes not dividing $30N\lvert d_K\rvert$, and $k\ge0$ an integer with $2c_5(\log r_1+\log r_2)\le16^k$. Let $L$ be a finite set of primes and, for $l\in L$, let $O_l$ be a coefficient ring as in Definition~\ref{w4:def:generic} and $X_l$ a set of Dirichlet characters modulo $N$ with values in $O_l$, each in generic position for $(K,l,e)$. Then every $\Sha_l(\Q,e;\psi)$ with $l\in L$ and $\psi\in X_l$ is finite, and
\[
\sum_{l\in L}\log l\sum_{\psi\in X_l}\frac{\len_{O_l}\Sha_l(\Q,e;\psi)}{e(O_l)}\le2c_5(\log r_1+\log r_2)\,N\,\bigl(N(e+1)\bigr)^A\le(Nt)^{k+1+A}.
\]
\end{proposition}

\begin{proof}
For a subgroup $H'\subset(\Z/N)^\times$ let $\mathcal O_l(H')$ be the set of characters with values in $O_l$ whose kernel is $H'$. Let $T$ be the set of kernels of the characters in $\bigcup_lX_l$. Each $X_l$ is contained in $\bigcup_{H'\in T}\mathcal O_l(H')$, and for $H'\in T$ let $L(H')$ be the set of $l\in L$ such that $X_l$ meets $\mathcal O_l(H')$. For $l\in L(H')$ some character of $\mathcal O_l(H')$ is in generic position, so (CorGK) with $H=H'$, $p=r_1$, $p'=r_2$ and the set $L(H')$ gives finiteness and
\[
\sum_{l\in L(H')}\log l\sum_{\psi\in\mathcal O_l(H')}\frac{\len\Sha_l(\Q,e;\psi)}{e(O_l)}\le2c_5\,\varphi([(\Z/N)^\times:H'])\,\bigl(N(e+1)\bigr)^A(\log r_1+\log r_2).
\]
All terms are nonnegative, so summing over $H'\in T$ bounds the left side of the proposition. The quotient $(\Z/N)^\times/H'$ is cyclic for $H'\in T$, so exactly $\varphi([(\Z/N)^\times:H'])$ complex characters of $(\Z/N)^\times$ have kernel $H'$. Distinct $H'$ give disjoint sets of characters, hence $\sum_{H'\in T}\varphi([(\Z/N)^\times:H'])\le\varphi(N)\le N$. This gives the first inequality. For the second, $2c_5(\log r_1+\log r_2)\le16^k\le N^k$, $N\le Nt$ and $N(e+1)\le Nt$.
\end{proof}

The hypotheses $r_1,r_2\nmid30N\lvert d_K\rvert$ and the generic position of the characters are the points at which the host enters. \notgiven{The verification that the characters of the host attached to a twist $e$ and a prime $l$ outside the finite set of bad primes with $(l-1)\nmid2e$ are in generic position for $K=\Q(\sqrt{-p_K})$, and that the host level $N$ is polynomial in the conductor of $E$, are not written out here.}

\subsection{A consequence in the direction of Vandiver's conjecture}\label{tw:sec:vandiver}

Consider the following statement.
\begin{enumerate}
\item[(V)] There are constants $C,A$ such that for every even $e\ge2$ and every finite set $L$ of primes $l\ge5$ with $(l-1)\nmid2e$, each $\Sha_l(\mathbb Q,e)$ is finite and $\sum_{l\in L}\lambda\bigl(\Sha_l(\mathbb Q,e)\bigr)\le C(e+1)^A$.
\end{enumerate}
The formal proof contains (V) as a theorem, derived from the statement that also gives the generic plus part of (I).

\begin{corollary}[Vanishing of eigenspaces]\label{w4:cor:large}
If (V) holds with constants $C,A$, then for every even $e\ge2$ and every prime $l\ge5$ with $(l-1)\nmid2e$ and $\log l>C(e+1)^A$, the $\omega^{-e}$ eigenspace of the $l$-class group of $\Q(\mu_l)$ vanishes.
\end{corollary}

\begin{proof}
If it did not vanish, $\Sha_l(\Q,e)\ne0$ by Proposition~\ref{w4:prop:classgroup}, so $\lambda(\Sha_l(\Q,e))\ge\log l$, and (V) with $L=\{l\}$ gives $\log l\le C(e+1)^A$.
\end{proof}

The exclusion $(l-1)\mid2e$ removes nothing of substance. If $(l-1)\mid2e$ and $(l-1)\nmid e$, then $\omega^{-e}$ is the quadratic character of $\Delta$, whose field is $\mathbb Q(\sqrt l)$, and the corresponding eigenspace is the $l$-part of a class group whose order is smaller than $l$.

\paragraph{Comparison with published results.} Kurihara~\cite{kurihara92} proved $A^{(\omega^{-2})}=0$ for every $l$, through $K_4(\mathbb Z)=0$, which is the case $e=2$ of Corollary~\ref{w4:cor:large} with no threshold. With $n=e+1$ odd, Soul\'e~\cite{soule99} proved $A^{(\omega^{-e})}=0$ once $\log l>n^{224n^4}$, a threshold of size $\exp(O(e^4\log e))$. De~la~Torre~Dur\'an~\cite{delatorre} improved the exponential factor for $F=\mathbb Q$ from $e^{O(n^4\log n)}$ to $e^{O(n^2\log n)}$, a threshold of size $\exp(O(e^2\log e))$. Corollary~\ref{w4:cor:large} has the threshold $C(e+1)^A$, one exponential smaller. In terms of $T(e)=\sum_l\lambda(\Sha_l(\mathbb Q,e))$, the published methods bound $\log T(e)$ by a constant times $e^2\log e$ at best, and (V) asserts $\log T(e)\le A\log(e+1)+\log C$. The methods of Soul\'e and of de~la~Torre~Dur\'an count cells of the Voronoi complex of $\mathrm{GL}_n(\mathbb Z)$, and the number of perfect forms grows like $e^{\Theta(n^2\log n)}$~\cite{vanwoerden}, so cell counting alone does not reach a polynomial bound. No finite computation can contradict (V), because its constants are existential. The computed groups $K_4(\mathbb Z)=0$ and $K_8(\mathbb Z)=0$ and the verification of Vandiver's conjecture for $l<2^{31}$ are consistent with every upper bound, and they give no information about (V).

\begin{remark}[odd twists]\label{w4:rem:odd}
For odd $e$ the analogous sum is classically polynomial. By the identification with $K_{2e}(\mathbb Z)\otimes\Zl$ in Subsection~\ref{tw:sec:sha} and the computation of $K_{4k-2}(\Z)$ from the Main Conjecture of Mazur and Wiles (Weibel's survey), $|\Sha_l(\Q,e)|$ is the $l$-part of the numerator of $B_{e+1}/(e+1)$ for $e=2k-1$ and $(l-1)\nmid e$. Hence $\sum_l\lambda(\Sha_l(\Q,e))\le\log|B_{e+1}|+o(e)=(1+o(1))\,e\log e$, using that the denominator of $B_{e+1}/(e+1)$ has logarithm $o(e)$. So the shape of (V) is the natural even analogue of a known odd statement. This supports plausibility and proves nothing.
\end{remark}

\paragraph{The chain of statements.} Let $K$ be an imaginary quadratic field. For a Dirichlet character $\psi$ write $\Sha_l(\mathbb Q,e;\psi)$ for the group of Definition~\ref{w4:def:sha} for the character $\psi^{-1}\chi^{-e}$ in place of $\chi^{-e}$, with coefficients in a finite extension $O$ of $\Zl$ that contains the values of $\psi$.

\begin{definition}[generic position]\label{w4:def:generic}
Let $K$ be a number field, $l$ a prime, $e\ge0$ and $N\ge1$ integers, $O$ a finite extension of $\Zl$ which is a discrete valuation ring, and $\psi$ a Dirichlet character modulo $N$ with values in $O$. For a prime $r$ let $\psi_{(r)}$ be the restriction of $\psi$ to the factor $(\Z/r^{v_r(N)})^\times$ of $(\Z/N)^\times$, viewed again as a character modulo $N$, and put $\psi^\circ=\psi\,\psi_{(l)}^{-1}$. Let $\omega\colon G_\Q\to\mathbb F_l^\times$ be the reduction of the cyclotomic character modulo $l$. We say that $\psi$ is in generic position for $(K,l,e)$ when the following five conditions hold.
\begin{enumerate}[leftmargin=*,label=(\alph*)]
\item $l\ne2$, $e\ge1$ and $N$ is even.
\item $l$ does not divide the order of $\psi$, $\psi_{(l)}^2=1$, and $\psi(-1)=(-1)^e$.
\item $K$ is imaginary quadratic, $l\nmid\lvert d_K\rvert h_K$, and $\gcd(d_K,N)=1$.
\item The subgroup $\ker\omega\cap\ker\psi^\circ$ of $G_\Q$ is not contained in $G_K$.
\item $(l-1)\nmid2e$.
\end{enumerate}
\end{definition}

For $\psi=1$ one has $\psi_{(l)}=\psi^\circ=1$, and conditions (b) and (d) become the conditions written in the proof of Proposition~\ref{w4:prop:corToBound}. For a general $\psi$ the condition in (b) concerns the $l$-component $\psi_{(l)}$ and not $\psi$ itself, and the condition in (d) concerns $\psi^\circ$.

For $\psi$ as above and $e\ge1$ we write $M_{e,O}(\psi^{-1})$ for $O\otimes_{\Zl}\Ql/\Zl$ with $\sigma\in G_\Q$ acting by multiplication by $\psi^{-1}(\sigma)\chi(\sigma)^{-e}$, where $\psi$ is viewed as a character of $G_\Q$ through $\Gal(\Q(\mu_N)/\Q)\cong(\Z/N)^\times$. Then $\Sha_l(\Q,e;\psi)$ is the subgroup of $H^1(G_\Q,M_{e,O}(\psi^{-1}))$ of the classes that vanish on $D_l$ and on $I_v$ for every finite $v$.

\begin{lemma}\label{w4:lem:unramchi}
Let $v\ne l$ be a prime, $e\ge1$, and let $\psi$ have order prime to $l$. A class in $H^1(D_v,M_{e,O}(\psi^{-1}))$ that vanishes on $I_v$ is zero. In particular $\Sha_l(\Q,e;\psi)$ is also the group of classes that vanish on $D_v$ for every $v\mid lN$ and on $I_v$ for every finite $v$.
\end{lemma}

\begin{proof}
Write $M=M_{e,O}(\psi^{-1})$ and let $\mathfrak m$ be the maximal ideal of $O$. Inflation and restriction give the exact sequence $0\to H^1(D_v/I_v,M^{I_v})\to H^1(D_v,M)\to H^1(I_v,M)$. Suppose first that $\psi$ is ramified at $v$, and choose $g\in I_v$ with $\psi(g)\ne1$. The element $\psi(g)$ is a root of unity of order $m>1$ prime to $l$, so $1-\psi(g)$ divides $m$ in $O[\psi(g)]$ and $1-\psi(g)\in O^\times$. The group $I_v$ acts trivially through $\chi$ because $v\ne l$, so $g$ acts on $M$ by $\psi(g)^{-1}$, and $(\psi(g)^{-1}-1)m'=0$ forces $m'=0$. Hence $M^{I_v}=0$. Suppose next that $\psi$ is unramified at $v$. Then $M^{I_v}=M$, and a Frobenius element $\varphi$ acts by $a=\psi(v)^{-1}v^{-e}$. If $a=1$, then $v^e$ would be a root of unity in the field of fractions of $O$, which is impossible because $v^e>1$. So $a-1\ne0$, multiplication by $a-1$ is surjective on the divisible $O$-module $M$, and $H^1(D_v/I_v,M)\cong M/(\varphi-1)M=0$. In both cases the inflated group vanishes. For the last assertion, a class that vanishes on $I_v$ for $v\mid N$, $v\ne l$, then vanishes on $D_v$, and conversely vanishing on $D_v$ implies vanishing on $I_v$.
\end{proof}

\begin{proposition}[(V) from (CorGK)]\label{w4:prop:corToBound}
Assume the following statement (CorGK) for $K=\Q(\sqrt{-3})$. There are $c_5,A$ such that for every even $N\ge2$, every subgroup $H\subset(\Z/N)^\times$, all distinct primes $p,p'\nmid30N\lvert d_K\rvert$, every finite set $L$ of primes such that for each $l\in L$ some character $\psi$ in the orbit of characters with kernel $H$ is in generic position in the sense of Definition~\ref{w4:def:generic}, one has finiteness and
\[
\sum_{l\in L}\log l\sum_{\psi}\frac{\len\Sha_l(\Q,e;\psi)}{e(O)}\le2c_5\,\varphi([(\Z/N)^\times:H])\,\bigl(N(e+1)\bigr)^A\bigl(\log p+\log p'\bigr),
\]
where $\Sha_l(\Q,e;\psi)$ is the group of Definition~\ref{w4:def:sha} for the character $\psi^{-1}\chi^{-e}$. Then (V) holds with $C=36c_52^A$ and the same $A$.
\end{proposition}

\begin{proof}
Take $N=2$, $H=(\Z/2)^\times$ (the trivial group), so the only character is $\psi=1$ and $\varphi(1)=1$, and take $p=7$, $p'=11$, which do not divide $30\cdot2\cdot3$. For $\psi=1$, $\Sha_l(\Q,e;1)=\Sha_l(\Q,e)$. The generic-position conditions for $K=\Q(\sqrt{-3})$, $\psi=1$, $l\ge5$, $e\ge2$ even and $(l-1)\nmid2e$ hold, as follows. (a) $l$ odd, $e\ge1$, $N=2$ even. (b) $l$ does not divide the order $1$ of $\psi$, $\psi^2=1$, and $\psi(-1)=1=(-1)^e$. (c) $K$ is imaginary quadratic with discriminant $-3$ and class number $1$, so $l\nmid3$; and $\gcd(3,2)=1$. (d) The last condition asks that $\ker(\chi\bmod l)\cap\ker\psi$ is not contained in $G_K$. Here $\ker\psi=G_\Q$, so it asks that $K\not\subset\Q(\mu_l)$, which holds because $3$ ramifies in $K$ and not in $\Q(\mu_l)$ for $l\ne3$. (e) $(l-1)\nmid2e$. The right side becomes $2c_5(2(e+1))^A(\log7+\log11)\le2c_5\,2^A(e+1)^A\cdot 18$. Lemma~\ref{w4:lem:coef} converts the left side into the left side of (V).
\end{proof}

\begin{proposition}[(CorGK) from (PropGK)]\label{w4:prop:prop2cor}
Assume the following statement (PropGK). For a number field $K$ there are $c_5,A$ such that for $e,N,H$ and every prime $p\nmid30N\lvert d_K\rvert$ there is an integer $NN_p\ne0$ with $\log\lvert NN_p\rvert\le c_5\varphi(\mathrm{index}\,H)(N(e+1))^A\log p$, and for every prime $l\ne p$ and every coefficient ring at which some character of the orbit is in generic position, the lengths are finite and $\sum_\psi\len\Sha_l(\Q,e;\psi)/e(O)\le2v_l(NN_p)$. Then (CorGK) holds with the same $c_5,A$.
\end{proposition}

\begin{proof}
Fix distinct admissible $p,p'$ and a finite set $L$ of primes in generic position. For $l\ne p$ use $NN_p$, and for $l=p$ use $NN_{p'}$. Then
\[
\sum_{l\in L}\log l\sum_\psi\frac{\len\Sha_l(\Q,e;\psi)}{e(O)}\le2\sum_l\log l\,\bigl(v_l(NN_p)+v_l(NN_{p'})\bigr)\le2\bigl(\log\lvert NN_p\rvert+\log\lvert NN_{p'}\rvert\bigr),
\]
because $\sum_lv_l(n)\log l\le\log|n|$ for a nonzero integer $n$. Insert the bounds for $\log|NN_p|$ and $\log|NN_{p'}|$.
\end{proof}

\begin{proposition}[(PropGK) from (Count) and (OneChar)]\label{w4:prop:onechar2prop}
Let $K$ be an imaginary quadratic field, $p\ge7$ a prime with $p\nmid N\,d_K$, $w\mid p$ a place of $K$ and $\sigma$ a geometric Frobenius element at $w$. Let $\mathcal O$ be a finite set of roots of unity $z$ stable under $\operatorname{Aut}(\mathbb C)$, and put $t(z)=Nw^{-e}z+1+Nw^{e}z^{-1}$, the trace of $\sigma$ on the Eisenstein triple $\Psi\oplus1\oplus\Psi^{-1}$ suitably normalized. Assume the following two statements.
\begin{enumerate}[leftmargin=*,label=(\alph*)]
\item \emph{(Count.)} There is $NN\in\Z\setminus\{0\}$ with $\log|NN|\le c\,|\mathcal O|\,(N(e+1))^A\log p$, such that for every prime $l\ne2,p$, every $\iota\colon\overline{\Q}_l\cong\mathbb C$ and every choice of finite sets $X_z$ of pairwise non-isomorphic admissible $\rho\colon G_K\to GL_3(\overline{\Q}_l)$ (level $N$, weight $e$), one has $\operatorname{tr}\rho(\sigma)\ne t(z)$ for $\rho\in X_z$ and $\lvert NN\rvert_l\le\prod_z\prod_{\rho\in X_z}\lvert\operatorname{tr}\rho(\sigma)-t(z)\rvert_l$.
\item \emph{(OneChar.)} For each character $\psi$ in generic position, if no admissible $\rho$ has $\operatorname{tr}\rho(\sigma)=t$, then $\Sha_l(\Q,e;\psi)$ is finite and there is a finite set $X_\psi$ of pairwise non-isomorphic admissible $\rho$ with $\len\Sha_l(\Q,e;\psi)/e(O)\le2\sum_{\rho\in X_\psi}v_l(\operatorname{tr}\rho(\sigma)-t_\psi)$.
\end{enumerate}
Then the conclusion of (PropGK) holds at $p$ with $NN_p=NN$.
\end{proposition}

\begin{proof}
The orbit of characters $\psi$ with kernel $H$ is in bijection with an $\operatorname{Aut}(\mathbb C)$-stable set $\mathcal O$ of roots of unity, $z_\psi=\psi(Nw)^{-1}$ under $\iota$, and $t_\psi=t(z_\psi)$; the number of elements is $\varphi(\mathrm{index}\,H)$. By~(a) no admissible $\rho$ has trace $t_\psi$, so~(b) applies to each $\psi$. Summing~(b) over $\psi$ and applying~(a) to the sets $X_\psi$,
\[
\sum_\psi\frac{\len\Sha_l(\Q,e;\psi)}{e(O)}\le2\sum_\psi\sum_{\rho\in X_\psi}v_l\bigl(\operatorname{tr}\rho(\sigma)-t_\psi\bigr)\le2v_l(NN).\qedhere
\]
\end{proof}

The content of (Count) is classical in type. The admissible $\rho$ correspond to Frobenius eigenvalues at $w$ of algebraic automorphic forms on a definite unitary group in three variables of level $N$ and weight of height $e$, whose number is polynomial in $N(e+1)$. The number $NN$ is the product of the scaled differences, rational by Galois stability, and nonzero because purity gives $|\operatorname{tr}\rho(\sigma)|\le3$ while $|t(z)|\ge p^{e}-2>3$. \notgiven{Statement (Count) is not proved here.}

\begin{proposition}[a consequence of the Selmer bound]\label{w4:prop:need}
Assume Hypothesis~\ref{w4:thm:vwin} for one primitive character $\psi_0$ of conductor $N_0$, and assume that for every $\varepsilon\in(0,\frac1{52})$ and $J\ge1$ there are arbitrarily large $X$ for which a window pair $(q,d)$ at scale $X^{1/J}$ exists (Definition~\ref{w4:def:winpair}). Then for every $\varepsilon\in(0,\frac1{52})$ there are arbitrarily large real $x$ and primes $q_x,d_x$ such that for every finite set $L$ of odd primes $l\nmid N_0q_xd_x$,
\[
\sum_{l\in L}\ \sum_{\substack{1\le e<x^{401}\\(l-1)\nmid e}}\lambda\bigl(\Sha_l(\Q,e)\bigr)\le\exp\bigl(\tfrac12x^{\varepsilon}\bigr).
\]
Consequently the product $P(T)=\prod_{e<T}\prod_{l}\lvert\Sha_l(\Q,e)\rvert$ over odd $l\nmid N_0$ with $(l-1)\nmid e$, outside two primes that depend on $T$, satisfies $\log\log P(T)\le\frac12T^{\varepsilon/401}$ for arbitrarily large $T$, for every $\varepsilon>0$.
\end{proposition}

\begin{proof}
Choose any function $e(x)$ with $x^{2+8\varepsilon}\le2^{e(x)}\le2x^{2+8\varepsilon}$. Let $J,X_0$ be given by Hypothesis~\ref{w4:thm:vwin}, take $X\ge X_0$ with a window pair $(q,d)$ at scale $X^{1/J}$, and take the $Y$ that Hypothesis~\ref{w4:thm:vwin} provides. Put $x=Y$, which tends to infinity with $X$. Let $N=\operatorname{lcm}(2^{e(x)+2},N_0,q,d)$ and let $E$ be the fixed field in $\Q(\mu_N)$ of the subgroup of Definition~\ref{w4:def:chifield}. Hypothesis~\ref{w4:thm:vwin} gives $\sum_{l\in L}\sum_{1\le e<t}\lambda(\Sha_l(E,e))\le\exp(\frac12x^\varepsilon)$ with $t=\lfloor x^{401}\rfloor$. For odd $l\nmid N_0qd$ we have $l\nmid N$, and Proposition~\ref{w4:prop:res} gives $\lambda(\Sha_l(\Q,e))\le\lambda(\Sha_l(E,e))$ whenever $(l-1)\nmid e$. The terms are nonnegative, so the inequality follows. For the consequence take $T=x^{401}$ and logarithms.
\end{proof}

By the identification with $K_{2e}(\mathbb Z)\otimes\mathbb Z_l$ and the comparison above, the published estimates give at best $\log\log P(T)\ll T^{2}\log T$, and nothing better along any sequence of $T$. Proposition~\ref{w4:prop:need} needs $\log\log P(T)\le T^{\delta}$ along a sequence of $T$, for every $\delta>0$. So the restriction to windows does not reduce what Hypothesis~\ref{w4:thm:vwin} requires to a published bound. The Selmer bound on windows implies a bound in the direction of Vandiver's conjecture that is stronger than the published bounds of Soul\'e and of de~la~Torre~Dur\'an. Alp\"oge's announcement~\cite{note} describes this bound as a weakening of Vandiver's conjecture.

\paragraph{The statement (OneChar).} Fix $K$, $l\ge5$, $e$ and a character $\psi$ in generic position, put $\Psi=\chi^e\psi$, and consider the ordinary deformation problem for representations of $G_{\mathbb Q}$ into the group of unitary type in three variables with residual representation $\Psi\oplus1\oplus\Psi^{-1}$. The formal proof derives (OneChar) from four inputs.
\begin{enumerate}[leftmargin=*,label=(G\arabic*)]
\item \emph{Tangent injection.} A class of $\Sha_l(\mathbb Q,e;\psi)$ gives a tangent vector of weight $-1$ at the Eisenstein point of a graded model of the ordinary deformation ring, with an action of a torus (the argument of Ribet).
\item \emph{An intersection inequality.} On a trigraded Rees algebra, under hypotheses (M1) to (M5), the cotangent module of weight $-1$ modulo the differentials of the Selmer equations has length at most $2\sum_\xi v(\xi(a))$, the sum over representatives of the torus orbits of the two-signed points $\xi$, with $a=\operatorname{tr}r(\sigma)-t$. This is commutative algebra with no arithmetic input.
\item \emph{Finiteness and automorphy.} The ordinary deformation problem is finite over $O$ (hypothesis (M2)), and every ordinary point $r$ with values in a finite extension of $O$ that is irreducible on every open subgroup of the kernel of the similitude character restricts on $G_K$ to a representation $\rho$ such that $\rho\otimes\chi^{-1}$ is automorphic and conjugate self-dual.
\item \emph{Rigidity.} At each two-signed point the tangent space of the Selmer locus is one-dimensional (hypothesis (M4)), which is the vanishing of the adjoint Bloch--Kato Selmer group. At the reducible fixed points $W\oplus1$ (hypothesis (M3)) the Bloch--Kato Selmer groups of the twist $W$ of a newform of weight $2e+1$ and of $W\otimes\theta$ vanish.
\end{enumerate}
\notgiven{None of (G1) to (G4) is proved here, and (OneChar) is not proved here.}

The closest published theorem to (G3) is Theorem~1.1 of Allen, Newton and Thorne~\cite{ant20}. It allows any number of residual constituents, but its hypothesis~(4) requires each constituent to be conjugate self-dual, and its hypotheses (5) and (6c) require a place at which the seed is an unramified twist of the Steinberg representation. In the configuration of (G3), on $G_K$ the constituents $\Psi$ and $\Psi^{-1}$ are exchanged by conjugate duality, so they are not individually conjugate self-dual when $e\ne0$, and no Steinberg place is assumed. The earlier theorem of Thorne~\cite{thorne15} has stronger hypotheses. So (G3) is not a published theorem. In characteristic zero at irreducible points, (G4) has the shape of Theorem~A of Newton and Thorne~\cite{nt23}, which needs the image to be enormous and does not need residual irreducibility. For the cyclotomic image, the trace criterion in Subsection~\ref{tw:rankfour} applies once the image spans the matrix algebra and its regular semisimple elements are dense.

At a reducible point let $f$ have weight $k=2e+1$. The Frobenius polynomials identify $W$ with $V(f)(e)$ by Chebotarev density and Brauer--Nesbitt. Take a cyclotomic field $F$ on which $\theta$ is trivial. Since $1\le e\le k-1$ and $e\ne k/2$, Kato's Theorem~14.2(1)~\cite{kato04} gives $H^1_f(F,V(f)(e))=0$. Restriction to $F$ is injective in characteristic zero. It gives the two vanishings over $\mathbb Q$ in (G4).

For the rank-three crossing in (G3), write the odd relations of a minimal graded presentation as $My=0$. Multiplication by the adjugate gives $(\det M)y=0$. Thus on the open set where $\det M$ is a unit all the odd coordinates vanish. In the complementary dimension argument the section bound gives dimension at least $3d+2$, whereas the triangular locus has dimension less than $3d$, with $d$ the degree parameter of the presentation. Propagation still requires finiteness over the weight algebra, a point known to be automorphic on the initial component, and a crossing theorem along the successive thickenings at reducible primes.

\begin{remark}[normalizations in (Count) and (OneChar)]\label{w4:rem:normalization}
Let $\sigma_w$ be an arithmetic Frobenius element at $w$, so that $\chi(\sigma_w)=Nw$ and $\psi(\sigma_w)=\psi(Nw)$, and let $\sigma=\sigma_w^{-1}$ be the geometric Frobenius. With $\Psi=\chi^e\psi$ the trace of $\sigma$ on $\Psi\oplus1\oplus\Psi^{-1}$ is $Nw^{-e}\psi(Nw)^{-1}+1+Nw^{e}\psi(Nw)$. With $z_\psi=\psi(Nw)^{-1}$ this is
\[
t(z_\psi),\qquad t(z)=Nw^{-e}z+1+Nw^{e}z^{-1}.
\]
The involution $z\mapsto z^{-1}$ permutes $\mathcal O$. The admissible $\rho\colon G_K\to\mathrm{GL}_3(\overline{\Q}_l)$ in both statements are of two kinds. Either $\rho$ is irreducible, $\rho\otimes\chi^{-1}$ is automorphic and conjugate self-dual, and $\rho$ has level $N$ and Hodge--Tate weights $(e,0,-e)$, or $\rho$ is conjugate to $W|_{G_K}\oplus1$ for an irreducible $W\colon G_\Q\to\mathrm{GL}_2(\overline{\Q}_l)$ with $W\otimes\chi^{-e}$ attached to a newform of weight $2e+1$, $\det W$ of finite order and level $N\lvert d_K\rvert$. In both cases the Frobenius eigenvalues at $w$ have absolute value $1$ under $\iota$, so $\lvert\operatorname{tr}\rho(\sigma)\rvert\le3$, while $\lvert t(z)\rvert\ge Nw^{e}-1-Nw^{-e}>3$ for $p\ge7$ and $e\ge1$. The explicit bound $4c_5\varphi(\mathrm{index})\lvert d_K\rvert^5N^{18}e^7(\log p+\log p')$ quoted in Subsection~\ref{tw:sec:vwin} belongs to a variant of (CorGK) that is not used. The argument uses (CorGK) with an unspecified exponent $A$, through Proposition~\ref{w4:prop:genericplus}.
\end{remark}


\subsubsection{The rank-four counts}\label{tw:rankfour}

Cases (II) and (III) of the plus part use symplectic representations of
rank four. Both use automorphy on every irreducible component through the
residual point. If $A$ is the deformation ring, $J$ the automorphic ideal
and $\mathfrak p$ the prime of the point, the starting condition is
$J\subseteq\mathfrak q$ for one minimal prime $\mathfrak q\subseteq
\mathfrak p$. To obtain the condition for every such $\mathfrak q$ one
joins the components by a chain. At each adjacent pair one needs the
local crossing assertion. Induction along the chain then gives the desired
containment for the last component. Finiteness of the set of systems of
Hecke eigenvalues alone gives no such containment.

At a locally generic automorphic point the following observation verifies
the trace condition for enormity in the adjoint vanishing theorem
of Newton and Thorne~\cite{nt23}. Enlarge the coefficient field so that
the eigenvalues in question split. Let $H$ be the image on the cyclotomic
subgroup. Assume that $H$ spans $M_n$ and that its regular semisimple
elements are dense in its Zariski closure. Let $V\subset M_n$ be a nonzero
simple $H$-submodule under conjugation. For regular $h$ write
$h=\sum_\alpha\alpha e_{h,\alpha}$, with $e_{h,\alpha}$ the spectral
projectors. If $\operatorname{tr}(e_{h,\alpha}v)=0$ for every such $h$,
every $\alpha$ and every $v\in V$, then $\operatorname{tr}(hv)=0$.
By density this holds on $H$ and by linearity on its span. Nondegeneracy
of the trace pairing gives $V=0$, a contradiction. In the symplectic
case the projection from endomorphisms to the symplectic Lie algebra is
$X\mapsto(X-X^\dagger)/2$, with $X^\dagger$ the adjoint for the alternating
form. The finite local condition is the inertia coboundary condition away
from $l$ and the crystalline condition above $l$.

The two counts use the same rigidity bound. In case (III) the prime
$l$ lies outside the bad set, which contains $2$ and $3$, and $l-1\mid 2e$.
Thus $e\ge2$. In the trace separation estimate put $y=p^e$ at an auxiliary
prime $p\ge5$. The required strict inequality is
$7y/2+y+p+1<py$: indeed $(p-9/2)y\ge p^2/2>p+1$.
In the product over conjugate roots of unity of order $k$, one unit
modulo $k$ acts on all three roots at once. The orbit therefore contributes
at most $\varphi(k)$ choices. Repeated roots retain their multiplicities.

\paragraph{The host count.}
In case (V), $l$ belongs to the exceptional set $B_\chi$ and $e=1$. On the branch $R_1\le 2^{c_2(\psi)}$, let $O$ be the coefficient ring with ramification index $E_O$, and normalize $v_l(l)=1$. At distinct two-signed points $r_i$ write
\[
 S=\sum_i v_l\bigl(\operatorname{tr}r_i(\varphi)-\operatorname{tr}r_z(\varphi)\bigr),
\]
where $r_z$ is the Eisenstein representation and $\varphi$ is geometric Frobenius. The host deformation and intersection bounds give $\operatorname{length}_O V\le 2E_OS$ for the Selmer module under consideration. If all trace differences are nonzero and $S\log l\le (lNe)^C$, division by $E_O$ gives
\[
 E_O^{-1}\operatorname{length}_O V\log l\le 2(lNe)^C\le(lNe)^{C+1}.
\]
This deduction requires the host deformation bound and the polynomial estimate for both kinds of representations in Remark~\ref{w4:rem:normalization}.

\subsubsection{Automorphic points and base change}\label{tw:levelpoints}


Let $E/F$ be the CM extension in the level construction, unramified at every finite place. Then $[F:\mathbb Q]\ge2$, since a quadratic field over $\mathbb Q$ ramifies at a finite prime. The definite unitary quotient is compact. Its irreducible summand $\pi$ therefore occurs discretely, with finite-dimensional archimedean type. Labesse's base change (Corollaire~5.3 of \emph{Changement de base CM et s\'eries discr\`etes}) gives one isobaric sum of discrete constituents $\Pi_i$ for $\pi$. At split places it gives the split local component and at spherical places it gives stable Satake transfer. At a nonsplit unramified place with residue cardinality $q$ the upstairs cardinality is $q^2$.

By the residual spectrum theorem of M{\oe}glin and Waldspurger, each $\Pi_i$ is a Speh representation of length $d_i$ on a normalized cuspidal representation of rank $a_i$, with $\sum_i a_id_i=n$. Conjugate self-duality follows for its seed from uniqueness of cuspidal support for the conjugate self-dual constituent. In the normalized induction the exponents are $(d_i-1)/2-j$, $0\le j<d_i$. Reversing their order gives the subrepresentation realization; in the unnormalized constant term the modulus contributes $a_i((d_i-1)/2-j)$ to the reversed exponent. Fix the local factors of these same global constituents at every place. Each prescribed local component belongs to the corresponding induced product as a subquotient: restrict the tensor product to a finite set containing that place and every nonspherical factor, and take its Jordan--Hölder factors. At the remaining places where only a subquotient is prescribed, finite length gives one. The infinitesimal multisets add under induction. For highest weight $\lambda$ their union at infinity consists of $\lambda_j+(n-1)/2-j$, $0\le j<n$.

For each Speh block put $b_{ij}=d_i-2j-n$. Apply the Galois representation theorem for regular algebraic polarized cuspidal representations, in the local compatibility hypotheses of this section, to the cuspidal seeds twisted by $(b_{ij}+a_i-1)/2$. Take the direct sum of the resulting representations. It is continuous, semisimple and of rank $n$. Under normalized local Langlands, parabolic products give sums and absolute-value twists give the corresponding twists of the parameter. The Speh exponents followed by the common twist $(1-n)/2$ thus give precisely the parameter of this sum at split places. This calculation applies to every correspondence satisfying the product and twist laws. On a seed infinitesimal weight $z$ the Hodge--Tate weight is $-z-b_{ij}/2$. The local conditions above $l$ pass to the sum under the spherical or tame hypotheses of the attachment theorem. If the sum is irreducible it consists of a single summand.

Now fix a common Taylor--Wiles base, a permissible auxiliary level, the free module $M$ and its Hecke annihilator $J$. For $m\ne0$ choose a nonzero coordinate in a free basis. Density of admissible weights gives a specialization where that coordinate survives. Decompose into irreducible automorphic summands and project onto a nonzero isotypic component. Take the local factors and the preceding Galois representation on this same summand. The spherical Hecke algebra acts there by its Satake character, so the Hecke polynomial is the characteristic polynomial at inverse Frobenius. The determinant and local types follow on the same summand by the product and twist laws. Projection detects the specialized vector. Varying $m$ gives separating classical points at the chosen level, with common base and annihilator $J$.

In rank two over a totally real field, take a regular algebraic polarized
Hilbert modular representation in the common-parity normalization $(k,w)$.
Skinner~\cite[eq.~(1), Thm.~1]{skinner09} gives the Frobenius-semisimplified
Weil--Deligne parameter with the twist $|\det|^{-1/2}$ and Hodge degrees
$(w-k_\tau)/2$ and $(w+k_\tau-2)/2$. Their sum is $w-1$.
At a place above $l$ an Iwahori-fixed vector gives trivial inertia of the
parameter and hence semistability. In the spherical case the monodromy
also vanishes, giving crystallinity. Away from $l$, purity includes the
monodromy filtration~\cite[Cor.~7]{blasius}. Use inverse Frobenius for the
geometric eigenvalues. If an algebraic character $\eta$ was used to put
$\pi$ in this normalization, the representation for $\pi$ is
$r_{\pi\otimes\eta}\otimes r_\eta^{-1}$. These assertions concern the
rank-two Hilbert modular input.

\subsubsection{The effective cohomology bound}\label{tw:effectiveK}

Let $p$ be a prime, $L$ a totally imaginary number field of degree $d$ and
absolute discriminant $D$, and $S$ the set of places above $p$.
Suppose $L$ contains a primitive $p$th root of unity if $p$ is odd
and a primitive fourth root if $p=2$. Write
\[
 B_C(e,d,D)=\exp((2ed)^C\log(2D)).
\]
The bound required by the arithmetic window argument is uniform: for one
real $C_K$, all such $p,L$ and all integers $e\ge1$, the group
$A=H^2(G_{L,S},\mathbb Z_p(e+1))$ is finite and
$\log |A|\le B_{C_K}(e,d,D)$. In forming the cochains, fix an algebraic
closure, a prolongation of each $v\in S$ to it, and the cyclotomic
character on the quotient unramified outside $S$. The estimate holds for
each such datum.

\paragraph{The localization kernel.}
For the discrete module $M=(\mathbb Q_p/\mathbb Z_p)(-e)$ let
$s_L$ denote the length over $\mathbb Z_p$ of the kernel of
\[
 H^1(G_{L,S},M)\longrightarrow \prod_{v\in S} H^1(G_{L_v},M).
\]
Lengths can initially be infinite. The source identifies this length with
the invariant defined using an absolute Galois group and a fixed algebraic
closure. At the end of the argument one returns from this invariant to
the same local conditions. The two twists involved here are $-e$ in degree
one and $e+1$ in degree two.

\paragraph{Counting in cyclotomic layers.}
Put $k=e-1$. The auxiliary field construction supplies a totally imaginary
field $L'$ with $d'=[L':\mathbb Q]=m d$, maximal cyclotomic depth $c$
($\mu_{p^c}\subset L'$, $\mu_{p^{c+1}}\not\subset L'$), and
\[
 2e+1<\min(p^c,r_2(L')), \quad
 m\le(d+1)(4e+4), \quad p^c\le 2d',\quad
 D'\le D^m(p^c)^{d'},\quad s_L\le s_{L'}.
\]
Here $D'=|\Delta_{L'}|$, $r_2$ counts pairs of complex embeddings, and
$c\ge1$, with $c\ge2$ at $p=2$. The strict inequalities are used before
applying the odd-degree homology estimate.

Indeed, if $c_0$ is the cyclotomic depth of $L$, take the layer
$L'=L(\mu_{p^{c_0+h}})$ with $h$ the least integer for which
$p^h\ge4e+4$. The cyclotomic layer theorem gives $m=p^h$.
Thus $m\le p(4e+4)$, $c\ge c_0+h$, and $r_2(L')=md/2>2e+1$.
The cyclotomic degree over $\mathbb Q$ gives
$p\le d+1$ and $p^c\le 2d'$. To bound the different, let $z$ generate
the layer as a primitive $N$th root, $N=p^{c_0+h}$, and let $P$ be
its relative minimal polynomial. Writing $X^N-1=PQ$ and differentiating
at $z$ gives $N=zP'(z)Q(z)$. The different contains $P'(z)$ and hence
$N$. Its norm divides $N^{d'}$, giving the displayed discriminant bound.
For the length comparison use the tower descent description at levels
$0$ and $h$. These lengths are those of the cokernels of $T-1$ and
$T^{p^h}-1$ for the same twisted operator $T$. The second image lies in
the first by the geometric-sum identity, giving a surjection of
cokernels in the required direction.

In a cyclic Galois layer $F/L'$ let $\zeta$ have order $q=p^\nu$,
$\nu\ge c$, and fix a generator $\gamma$ of order $p^{\nu-c}$
with $\gamma\zeta=\zeta^{1+p^c}$. Set $a=(1+p^c)^e$. Count the
functions $\varphi$ from nonzero prime ideals of $\mathcal O_F$ to
$\mathbb Z/q$ satisfying
\[
 \sum_P\operatorname{ord}_P(\alpha)\varphi(P)=0\quad(\alpha\in F^\times),
 \qquad \varphi(\gamma^{-1}P)=a\varphi(P).
\]
The first condition makes these characters of the ideal class group.
Write $h_\nu$ for their number, $T_j$ for the order of the torsion
subgroup of $H_j(\mathrm{GL}_{4k+10}(\mathcal O_{L'}),\mathbb Z)$, and
\begin{align*}
 b&=2k+2+(2k+3)(c+2k+3), & f&=v_p(e),\\
 g&=2d'+\lfloor\log_2 h_{L'}\rfloor+2, & X&=(b+f)g+c+f .
\end{align*}

Two counts enter. For the first count, consider a finite subset of
$F^\times$ whose elements have
$\operatorname{ord}_P(\alpha)=\mathfrak d(\gamma P)-a\mathfrak d(P)$
for a divisor $\mathfrak d$ depending on $\alpha$. The chosen elements
are distinguished modulo $q$th powers by the $p^b$th powers of their
twisted norms. In the archive convention the norm is
$\prod_{j=0}^{|\gamma|-1}\gamma^j(\alpha)^{a^j}$, with the field
automorphism acting on $\alpha$. Each element gives a cycle
in degree $2k+3$. One chooses the residue tests simultaneously away
from $p$ and the finitely many denominators of these cycles, at primes
where $\mathcal O_{L'}$ surjects onto the residue field. If two cycles
have the same scaled evaluations at every such prime, the ratio of their
powered twisted norms is a $q$th power in every tested residue field.
The power-detection lemma gives a $q$th power in $F^\times$.
The separation condition then identifies the two elements. The image
bound for the resulting primitive family of evaluations bounds the first
count by $q^{r_2(L')}T_{2k+3}T_{2k+2}$.

\paragraph{The cycle and its determinant.}\label{kone:cycle}
Write $A_q=\mathbb Z/q$, $C=\mathbb Z/q$ as an additive group,
$M=2k+3$, $\kappa=1+p^c$ and $a=\kappa^e$, so $e=k+1$.
In the bar complex of $C$ with coefficients in $A_q$ use
\[
 z=\sum_{j_1,\ldots,j_e\in C}[j_1|1|\cdots|j_e|1].
\]
This has degree $2e$. In its boundary the first face sums to zero
because its multiplicity is $q$. The remaining faces cancel in pairs
after shifting a summation variable by one. Put $d_t(x)=(x,tx)$,
$0\le t\le M$. The product-group homology in degree $2e$ with these
coefficients has a generating set of size at most $2e+1=M$ (use the
first projection onto $C$). More explicitly, constants, cyclic
shift minus identity and summation give an exact sequence
from the trivial module through two copies of the cyclic permutation
module to the trivial module. The two connecting sequences and
induction from the kernel bound the generator number in degree $d+2$
by that in degree $d$ plus the kernel-group bounds in degrees $d+1,d+2$.
For a cyclic kernel both bounds are one; degree zero needs one
generator. Hence degree $2e$ needs at most $2e+1$ generators.
The $M+1$ classes $(d_t)_*z$
therefore satisfy a relation $\sum_t n_t(d_t)_*z=\partial B$ with
some $n_{t_0}$ a unit. Indeed their subgroup is a finite abelian
$p$-group with at most $M$ cyclic factors. A redundant generator
modulo $p$ is redundant in the subgroup, by successive lifting
through the powers of $p$. This gives the unit coefficient.

Index the coordinates of $F^{C^2}$ by $j=(j_1,j_2)$.
On coordinate $j$ the diagonal representation $\rho$ of $C^2$
has character $(x,y)\mapsto\zeta^{j_1x+j_2y}$. For $t$ let
$s_t(j)=j_1+tj_2$. The lattice construction uses a finitely supported
integer function $\eta$ on $C^2$. Its sums on the fibers of $s_t$
vanish for $t\ne t_0$. At $t_0$ its pushforward is
\[
 \sum_{i=0}^{M}[X^i]P(X)\,[\kappa^i],\qquad
 P(X)=\prod_{j=0}^{M-1}(X-\kappa^j).
\]
Here $[x]$ denotes a basis vector in $\mathbb Z[C]$. To see the
interpolation, use the product of $[(0,1)]-[(t,0)]$ over $t\ne t_0$
in $\mathbb Z[C^2]$. Each unwanted pushforward kills a factor. At
$t_0$ the product equals a translation times
$\prod_{t\ne t_0}([t_0-t]-1)$. Every $t_0-t$ generates a subgroup
containing $p^c$ in $C$, since $0<|t_0-t|\le M<p^c$. Thus each factor
divides $[p^c]-1$. The desired pushforward is divisible by
$([p^c]-1)^M$: in polynomial coordinates apply
$\prod_{j<M}(U-\kappa^j)$ to $X$, where $Uf(X)=f(X^\kappa)$.
The operator $U-\kappa^j$ carries $(X^{p^c}-1)^j$ times any polynomial
into $(X^{p^c}-1)^{j+1}$ times a polynomial, by factoring
$X^{\kappa p^c}-1$ and reducing the quotient modulo $X^{p^c}-1$.
Lift the quotient in the group ring by $x\mapsto(x,0)$ and undo the
translation. This constructs $\eta$ with both specified properties.

Modify the standard lattice coordinate $j$ by the fractional ideal
with divisor $\eta(j)\mathfrak d$. For each $t$ construct an
$L'$-linear $g_t$ taking the standard lattice to this modified lattice
and commuting with $\rho d_t$. The determinant-one identifications on
blocks use the Steinitz relation for a direct sum of fractional ideals:
one can move the sum of two divisors to one coordinate, and iterate.
Scaling one distinguished coordinate realizes a prescribed determinant
when its divisor is the difference of the two total divisors.
For $t\ne t_0$ all block determinants can be taken as one. At $t_0$
the construction groups coordinates in Galois orbits and leaves the
determinant
\[
 \beta=\prod_i\gamma^{-i}(\alpha)^{[X^i]P_e(X)},\qquad
 P_e(X)=\prod_{0\le j<M,\ j\ne e}(X-\kappa^j)
\]
at character $1$. Indeed the divisor of $\gamma^{-i}\alpha$ at $P'$
is $\mathfrak d(\gamma^{i+1}P')-a\mathfrak d(\gamma^iP')$.
Multiplying this difference by the coefficients of $P_e$ gives
precisely the divisor prescribed by $P=(X-a)P_e$. This explains
both the inverse action in $\beta$ and the missing factor in $P_e$.

Let $H$ be the stabilizer of the standard lattice inside the localized
group containing all $g_t$ and $\rho(C^2)$. One base-field denominator
excludes the support of the divisors and the bad primes of the splitting
map. Set $\rho'=g_{t_0}^{-1}\rho g_{t_0}$ and $h_t=g_{t_0}^{-1}g_t$.
Both representations take values in $H$, and $h_t\in H$ since both
maps $g_t,g_{t_0}$ take the standard lattice to the same lattice.
For the bar prism $P_h$ one has
$\partial P_h+P_h\partial=(\mathrm{conj}_{h^{-1}})_*-1$.
Consequently the chain
\[
 Y=\rho'_*B-\rho_*B+\sum_t n_t P_{h_t}((\rho'd_t)_*z)
\]
is closed in $H$. In fact $h_t^{-1}\rho'd_th_t=\rho d_t$, so the
boundary of the sum is $\rho_*\partial B-\rho'_*\partial B$.
In the localized group $Y$ is homologous to
$\sum_t n_t P_{g_t}((\rho d_t)_*z)$. This follows by applying the
prism identity to $\rho_*B$ and the prism product identity to
$g_t=g_{t_0}h_t$ on each cycle. Multiply $Y$ by $n_{t_0}^{-1}$.
The standard lattice is finite projective over $\mathcal O_{L'}$.
Adding a projective complement with trivial action realizes this cycle
in some $\mathrm{GL}_N(\mathcal O_{L'})$.

At an allowed residue field $k_0$, restriction of scalars splits into
the conjugate blocks indexed by $\sigma\in\mathrm{Gal}(F/L')$.
For a character $v:k_0^\times\to A_q$, the finite-field functional
on a scalar-character block of weight $x$ evaluates a prism for $g$
as $x^e v(\det g)$. On the scalar line this is the product of $e$
cyclic carry cocycles with $v$: each consecutive pair in $z$ sums to
$x$, and insertion in the prism leaves $v(g)$. Transfer over projective
lines gives the additive functional on vector spaces. The line count
for direct sums uses $|k_0|=1$ in $A_q$, since the reduced $\zeta$ has
order $q$. Additivity kills the trivial complement. Blocks of determinant
one give zero. In the surviving block let $\sigma\zeta=\zeta^{\chi(\sigma)}$.
Summing the conjugate blocks gives
\[
 \sum_\sigma \chi(\sigma)^e v(\overline{\sigma\beta})
   =P_e(a)\,v\big(\overline{\operatorname{nrm}(\alpha)}\big).
\]
Indeed replacing $\alpha$ by $\gamma^{-i}\alpha$ multiplies the
weighted sum by $a^i$ after reindexing, and the sum for $\alpha$ is
the evaluation of its twisted norm. The unit coefficient cancels the
normalization of $Y$. For a cyclic residue unit group take $v$ with
kernel the $q$th powers. Each factor $\kappa^e-\kappa^j$ of $P_e(a)$
has valuation $c+v_p(|e-j|)\le c+M$, by lifting the exponent
(with $c\ge2$ at $p=2$). In $A_q$, $P_e(a)$ therefore divides
$p^{M(c+M)}$. Equality of two evaluations scaled by $p^{2e}$
implies that the ratio of their twisted norms raised to
$p^{2e+M(c+M)}=p^b$ is a $q$th power at this prime.

\paragraph{Detecting powers and bounding the image.}\label{kone:rank}
If $\xi\in F^\times$ passes all these residue tests, choose
$x$ with $x^q=\xi$ in an algebraic closure. For each automorphism
$g$ over $F$, the condition $\sigma x=gx$ describes an open coset.
Degree-one Chebotarev gives a Frobenius in this coset outside the finite
exclusion set and the primes dividing $q$. Its residue field has prime
order, so reduction from $\mathcal O_{L'}$ is surjective. In the residue
field $\xi$ has a $q$th root and all the $q$th roots of unity are
present. Hence the residue of $x$ is fixed by Frobenius. The quotient
$\sigma x/x$ is a $q$th root of unity congruent to one. If it were
different from one, its geometric sum of length $q$ would be zero,
giving $q=0$ in the residue field. Thus $gx=x$ for every $g$ and
$x\in F$. Applying this argument to the ratio proves injectivity
for the simultaneous evaluations of the finite set.

For the image bound put $m_1=2e$, and stabilize the finitely many
cycles to rank $n\ge2(m_1+1)+4$. An integral lift $Z$ of a
mod-$q$ cycle has $\partial Z=qV$. The homology class of $V$ is
$q$-torsion in degree $m_1$, independent of the lift. When this class
vanishes one can subtract $q$ times an integral chain from $Z$ to
make it a cycle. There are at most $T_{m_1}$ fibers of this obstruction
by homological stability. Differences in a fiber therefore reduce the
count to evaluations on integral cycles.

Write $\psi_p$ for the $p$-fold block sum on integral homology and
$i$ for stabilization to the same rank. Starting with $z_0=x$ put
$z_{t+1}=(\psi_p)_*z_t-p^{t+2}i_*z_t$. A primitive evaluation
$\phi$ is stable and additive on block sums. Hence
$\phi(z_{t+1})=p(1-p^{t+1})\phi(z_t)$. The factors $1-p^{t+1}$
are units modulo $q$. The scaled evaluations $p^{m_1}\phi(x)$ thus
factor through the integral span of the $z_{m_1}$. Its rank is at
most $r_2(L')$. To obtain this last inequality the source passes to
real cycles and pairs with one odd-degree functional at each infinite
place. A class in the common kernel, if nonzero, survives stabilization
and is separated by a closed cochain. The surjectivity theorem represents
that cochain on cycles by an invariant form. Its weight decomposition
consists of the place functionals and terms of weights $j=2,\ldots,m_1+1$.
The latter vanish on the displayed span: the factor at $t=j-2$ is
$p^j-p^j=0$. The former vanish by the kernel condition. This contradicts
separation. The real image thus has dimension at most the number of
complex places, giving the integral rank bound.

For any finite set of evaluations use the finitely generated subgroup
spanned by lifts of its elements. Its free part contributes at most
$q^{r_2(L')}$ values. Its torsion contributes at most $T_{m_1+1}$
by stability, also at the enlarged rank for $z_{m_1}$. Multiplication
by the number of obstruction fibers gives the claimed bound.

\paragraph{Transfer over residue-field lines.}\label{kone:lines}
Here is the transfer used in the cycle calculation of
Paragraph~\ref{kone:cycle}. Let $V$ be a vector space over the finite
residue field $k_0$, let $\rho:G\to\mathrm{GL}(V)$, and let $t$ be the
functional on bars of $k_0^\times$ which vanishes on boundaries.
Choose a representative on each line. This gives a function
$\lambda:V\setminus\{0\}\to k_0^\times$ with
$\lambda(sv)=s\lambda(v)$, by writing $v$ as a multiple of its
representative. Give the lines the right action
$L\cdot g=\rho(g)^{-1}L$. For $0\ne v\in L$, put
\[
 c(g,L)=\lambda(v)/\lambda(\rho(g)^{-1}v).
\]
The ratio is independent of $v$. It satisfies
$c(gh,L)=c(g,L)c(h,L\cdot g)$ by cancellation at the intermediate
line. For a bar $[g_1|\cdots|g_n]$ sum, over initial lines $L_0$,
the bars $[c(g_1,L_0)|\cdots|c(g_n,L_{n-1})]$, where
$L_i=L_{i-1}\cdot g_i$. This sum commutes with the bar boundary.
For the first face reindex by $L_1$; for an interior face use the
cocycle identity; the last face just forgets the last arrow.
Composition with $t$ therefore kills boundaries.

If $\lambda$ is multiplied by a scalar $b(L)$ on line $L$, each
edge label is multiplied by $b(L_{i-1})/b(L_i)$. The two labelled
paths are related by the vertical edges $b(L_i)$: the products along
the two sides of each square agree. Triangulate the prism on each
path by inserting the vertical edge successively at the $n+1$
vertices, with alternating signs. Its interior faces cancel in pairs;
the two ends give the difference of the old and new bars, and the
side faces are the corresponding prisms on the boundary. Thus the
two sums agree under $t$ on cycles. In particular a linear change
of coordinates, which transports both lines and representatives,
preserves the resulting functional.

On $V\oplus W$, send a line to its projection in $V$ when that
projection is nonzero, and otherwise to its line in $W$. Use the
normalizing function of the indicated projection. Each line above
a fixed line in $V$ is the graph of a linear map from that line to
$W$, so its fiber has $|k_0|^{\dim W}$ elements. A line in $W$ has
one preimage. For block diagonal actions the edge labels in a fiber
are exactly those of its projected path. Both fiber sizes equal one
as coefficients in $\mathbb Z/q$, since $q\mid |k_0|-1$. The
functional on a direct sum is consequently the sum of the two
functionals on cycles. On the one-dimensional scalar representation
there is a single line and $c(g,L)=g$, so the construction recovers
$t$. These two facts justify passage from the scalar computation to
the sum over the residue-field blocks.



For the second count write $Y_F=F^\times$ additively, let $I_F$ be the
divisor group, and let $\partial:Y_F\to I_F$ be the divisor map.
On elements put $\delta=\gamma^{-1}-a$; on divisors the corresponding
operator is $(\delta_I d)(P)=d(\gamma P)-a d(P)$. Put
\[
 U=\partial^{-1}(\delta_I I_F),\quad
 W=U/(\delta Y_F+Y_{F,\mathrm{tors}}),\quad
 C'=I_F/(\partial Y_F+\delta_I I_F).
\]
Let $E'$ be $\ker\partial$ modulo $\delta(\ker\partial)$ and torsion.
The divisor and unit counts give finite groups with $|W|=|E'||C'|$,
$h_\nu\le p^{v_p(|C'|)}$, and
\[
 v_p(|E'|)+c+f=r_2(L')(\nu+f).
\]
The last equality uses $v_p(a-1)=c+f$ and
$v_p(a^{p^{\nu-c}}-1)=\nu+f$.

Here are the two order computations. On $I_F$ the operator $\delta_I$
is injective, since an element in its kernel satisfies
$(a^{|\gamma|}-1)d=0$ and the divisor group is torsion free.
For $y\in U$ the equation $\partial y=\delta_I d$ now determines $d$.
Sending $y$ to the ideal class of $d$ maps $W$ onto the kernel of
$\delta_I$ on the class group. Indeed a class in that kernel has
$\delta_I d$ principal, and altering $y$ by $\delta z$ alters $d$ by
$\partial z$. If $d$ is principal one can therefore take $y$ to be
a unit. Two such units agree in $W$ precisely modulo the twisted
differences of units and torsion, again by injectivity on divisors.
Thus the kernel is $E'$. The kernel and cokernel of an endomorphism
of a finite group have the same cardinality, and here the cokernel
is $C'$. Equivariant class characters factor through $C'$; on each
cyclic $p$-summand of order $p^i$ there are at most $p^i$ choices.

The logarithmic unit lattice spans the real hyperplane of coordinates
at infinite places with sum zero. On all coordinates $\gamma^{-1}$
permutes $r_2(L')$ cycles of length $n=p^{\nu-c}$, since each complex
place downstairs has $n$ lifts. The absolute determinant of
$\gamma^{-1}-a$ on the permutation space is $(a^n-1)^{r_2(L')}$.
Its absolute determinant on the one-dimensional quotient is $a-1$.
The determinant on the lattice modulo roots of unity therefore gives
$|E'|(a-1)=(a^n-1)^{r_2(L')}$, as required.

The powered twisted norm induces a map
$W\to Y_F/(qY_F+Y_{F,\mathrm{tors}})$. Its image has size at most
$B=q^{r_2(L')}T_{2k+3}T_{2k+2}$ by choosing one representative for
each image element and applying the first count. Write
$a^{p^{\nu-c}}-1=p^{\nu+f}m_0$ with $p\nmid m_0$.
The norm identity kills the kernel by $p^{b+f}m_0$.
Indeed, write $g=\gamma^{-1}$, $n=|\gamma|$, $N_a=\sum_{j<n}a^jg^{n-j}$.
Then $(g-a)N_a=N_a(g-a)=(1-a^n)g$. If $p^bN_a y=qz+t$
with $t$ torsion, applying $g-a$ and dividing by $q$ modulo torsion
shows that $p^{b+f}m_0g y$ lies in $(g-a)Y_F+Y_{F,\mathrm{tors}}$.
This subgroup is invariant under $g$ and $g^{-1}$, proving the assertion.
On its $p$-primary subgroup this gives an exponent bound $p^{b+f}$.
The Selmer injection and the base-field Selmer count bound $|W[p]|$
by $p^g$. A finite abelian $p$-group killed by $p^{b+f}$, with
at most $p^g$ elements killed by $p$, has order at most $p^{(b+f)g}$.
Indeed multiplication by $p$ takes the subgroup killed by $p^{j+1}$
into the subgroup killed by $p^j$ with kernel of order at most $p^g$.
Induction from $j=0$ gives the bound. Applying it to the kernel gives
$p^{v_p(|W|)}\le Bp^{(b+f)g}$. In particular
\[
 h_\nu q^{r_2(L')}\le
 p^{v_p(|C'|)+v_p(|E'|)+c+f}
 =p^{v_p(|W|)+c+f}\le Bp^X.
\]
For the Selmer bound just used, lift $w\in W[p]$ to $u\in U$ and write
$pu=\delta y+t$ with $t$ torsion. If $\partial u=\delta_I d$, then
$\partial y=pd$ by injectivity of $\delta_I$. Also $gy-y$ vanishes
modulo $pY_F+$ torsion, since $p\mid a-1$. The class of the chosen
$y$ in that quotient distinguishes the $w$: equality of two such
classes gives $y-y'=pz+t'$, whence
$p(u-u'-\delta z)$ is torsion, as is $u-u'-\delta z$ itself.
To descend these classes, fix a generator $\rho$ of the roots of unity
of $F$. In multiplicative notation a representative $\alpha$ has
$\gamma\alpha/\alpha=v^p\rho^r$ for $0\le r<p$ after absorbing a
$p$th power. The pair $(r,N_{F/L'}v)$ takes at most $p^2$ values
because $(N_{F/L'}v)^p=(N_{F/L'}\rho)^{-r}$. Within a fiber the ratio
of two such $v$ has norm one. Hilbert 90 supplies a coboundary, so
after removal of a $p$th power the ratio of the two $\alpha$ descends
to $L'$. Outside $p$ the cyclotomic extension is unramified, so the
valuations of the descended element are divisible by $p$ there.
The resulting base-field classes modulo powers have size at most
$p^{2d'+\lfloor\log_2 h_{L'}\rfloor}$. Indeed valuations modulo $p$
at places over $p$ cost at most $p^{d'}$. With these fixed, taking
the $p$th root of the principal divisor costs at most
$|\mathrm{Cl}(L')[p]|\le p^{\lfloor\log_2 h_{L'}\rfloor}$.
The remaining units modulo $p$th powers cost at most $p^{d'}$
by the unit theorem. This gives the stated $p^g$.

Substituting the value of $B$ gives
\[
 h_\nu q^{r_2(L')}
 \ \le\ q^{r_2(L')}T_{2k+3}T_{2k+2}\,p^X.
\]
Cancel the positive common factor. All quantities on the right of
$h_\nu\le T_{2k+3}T_{2k+2}p^X$ are independent of $\nu$.
This uniformity allows passage back to the localization kernel:
restriction to the infinite cyclotomic tower is injective on that kernel,
and each finite set in its image is detected at a common finite level.
At that level the corresponding torsion class characters inject into
the functions counted above. Thus every such finite set has size at most
$T_{2k+3}T_{2k+2}p^X$. It follows that the kernel is finite and
\[
 s_L\log p\le s_{L'}\log p
 \le\log T_{2k+3}+\log T_{2k+2}+X\log p.
\]

For clarity on the constants, write $u=2ed'$ and $Y=\log(2D')$.
For the torsion estimate set $N=4e+6$ and $d_* = d'N$.
The integral face complex for $\mathrm{GL}_N(\mathcal O_{L'})$ is
exact in the range used to compute $H_j$. In dimensions up to $j$
its number of orbits is bounded by
$A_*=(2D'd_*)^{12(j+1)d_*^3N}$ and the orders of the stabilizers by
$s_*=3^{d_*N}$. The boundary from dimension $j+1$ is the alternating
sum of $j+2$ faces, hence has column $\ell^1$-norm at most $j+2$.
The equivariant chain estimate with these bounds gives
\[
 T_j\le (j+2)^{A_*}\,s_*^{\,jA_*s_*^j}.
\]
This estimate uses exactness of the augmented complex, orbit counting
for integral lattices, and the bounds for homology of finite stabilizers.
Taking logarithms gives
$A_*\log(j+2)+jA_*s_*^j\log s_*$. The numerical estimate for this
expression at $j\le2e+1$ gives
$\log T_j\le \exp(u^{C_0}Y)$ with $C_0=24$ in the two degrees above. The numerical
comparison for the layer gives $X\log p\le u^9Y$. Absorbing the three
terms and descending using the degree and discriminant bounds gives
$s_L\log p\le B_C(e,d,D)$ with $C=5(\max(C_0,9)+2)+11$.

\paragraph{From the localization kernel to the required order.}
Poitou--Tate duality and the local finiteness and duality inputs give
\[
 \operatorname{length}_{\mathbb Z_p} A
 \le s_L+\sum_{v\in S}\operatorname{length}_{\mathbb Z_p}H^0(G_{L_v},M).
\]
For each $v$ the cyclotomic image $U_v\subset\mathbb Z_p^\times$ has
index $I_v\le d$: it contains the inertia image, whose index is at most
$d$. The element $(1+p^2)^{I_v}$ lies in $U_v$.
An invariant of order $p^t$ forces
$p^t\mid (1+p^2)^{I_v e}-1$. The valuation of the latter is
$2+v_p(I_v)+v_p(e)$, also for $p=2$ since the base is $1\bmod 4$.
Consequently the invariants form a subgroup of a cyclic group of order
at most $p^{2+\lfloor\log_p d\rfloor+e}$. There are at most $d$ places
above $p$, so
\[
 \operatorname{length}_{\mathbb Z_p} A
 \le s_L+d(2+\lfloor\log_p d\rfloor+e).
\]
A finite-length $\mathbb Z_p$-module of length $t$ has order $p^t$,
by a composition series with successive quotients $\mathbb F_p$.
This proves finiteness. The root-of-unity hypothesis gives $p\le d+1$,
and $\lfloor\log_p d\rfloor\log p\le\log(d+1)$. Hence the extra term
in $\log|A|$ is at most $d(3+e)\log(d+1)$. For $e,d,D\ge1$ this is
at most $4d\log(2D)(2e)(2e+3)\log(d+1)$, which the final numerical
comparison absorbs into $B_{\max(C,0)+9}(e,d,D)$ together with $B_C$.
Following the exhibited witnesses gives $5(24+2)+11=141$ at the
localization-kernel step and $141+9=150$ here. The interface used for
the window only asks for the existence of one absolute exponent.

\paragraph{Scope of this reconstruction.}
This deduction uses the arithmetic and analytic inputs to the odd-degree
rank bound, the integral torsion estimate, the Selmer and class-group
comparisons, and the tower and duality theorems in their indicated
normalizations. Paragraph~\ref{kone:lines} gives the line transfer used in the cycle
argument. A standalone derivation still requires the integral lattice
and splitting results and the arithmetic and analytic inputs above. In the window argument
the same uniform bound is supplied to the case estimates together with the
separate rank-two Selmer input.

\paragraph{Use of the bound on the selected family.}\label{kone:consumer}
The window argument takes the existential pair $(C_K,\text{cohomology
bound at }C_K)$ once. It uses that same pair in the count of type I,
in the joint count of types II and III, and in the count of type V.
Type I also uses the rank-two Selmer estimate and the unramified
automorphy input. Types II and III use the symplectic input, while
type V uses the reduced unitary lifting input. The four counts and
the common pair yield the window estimate of
Hypothesis~\ref{w4:thm:vwin} in the source normalization.
At each selected index the window supplies finiteness of the odd
localization kernels and bounds the sum of their lengths, weighted by
$\log p$, by $\sqrt b$. Indeed the record identifies its finite
lengths with the cohomological ones at every $1\le e<t$; nonprime
indices contribute zero, and summation gives the required inequality.
With the local lattice comparison and equality of the small-level
ranks at that index this gives the twisted-place divisibility used
in the division estimate of Subsection~\ref{sec:frontier}.
These rank and automorphy hypotheses enter independently of the
cohomological order bound.



\subsection{The congruence bound}\label{tw:sec:linkn}

\begin{hypothesis}[congruence bound]\label{w4:thm:linkn}
There are $C$ and $x_0$ such that for every admissible datum at scale $x\ge x_0$ for which every $\Sha_l(E,e)$ with $0<e<t_r$ is finite, and for every lift index $j$,
\[
\sum_{w\ \mathrm{twisted}}\ell_we_w\,\eta_j(w)\ \le\ d^0_j\Bigl(C\bigl(f\,m\,b\,L_\varepsilon^2+(\log x)^2N_c\bigr)+\sum_l\log l\sum_{1\le e<t_r}\len\Sha_l(E,e)\Bigr),
\]
provided the last sum converges.
\end{hypothesis}

The argument for Hypothesis~\ref{w4:thm:linkn} has three steps. A lattice defect is bounded by congruence numbers (step one), congruence numbers by Fitting lengths of a cocycle module (step two), and Fitting lengths by Selmer lengths plus local terms (step three).

Let $O$ be a complete discrete valuation ring with fraction field $K$, $X$ a finite set, $V=\bigoplus_{x\in X}V_x$ a finite-dimensional $K$-vector space, and $\Lambda\subset V$ an $O$-lattice. Let $T\subset\prod_{x\in X}O$ be an $O$-subalgebra acting on $V$ by $t\cdot v=(t_xv_x)_x$ with $T\Lambda\subset\Lambda$. For $x\in X$ define the \emph{congruence ideal} $\mathfrak c_x=\{t_x: t\in T,\ t_y=0\text{ for all }y\ne x\}$ and $\delta_x=\min\{v(a):a\in\mathfrak c_x\}\in\mathbb N\cup\{\infty\}$. For $Z\subset X$ write $V_Z=\bigoplus_{x\in Z}V_x$ and $\pi_Z\colon V\to V_Z$ for the projection.

\begin{proposition}\label{w4:prop:cong}
For $Z\subset X$, put $M^{\rm up}=\pi_Z(\Lambda)$ and $M^{\rm low}=\Lambda\cap V_Z$. Then
\begin{enumerate}[leftmargin=*,label=(\roman*)]
\item $M^{\rm up}/M^{\rm low}\cong\Lambda/\bigl((\Lambda\cap V_Z)\oplus(\Lambda\cap V_{X\setminus Z})\bigr)$;
\item $\len_O(M^{\rm up}/M^{\rm low})\le\sum_{x\in Z}\dim_KV_x\cdot\delta_x$.
\end{enumerate}
\end{proposition}

\begin{proof}
(i) The map $\Lambda\to M^{\rm up}/M^{\rm low}$, $v\mapsto\pi_Z(v)$, is surjective. Its kernel is the set of $v\in\Lambda$ with $\pi_Z(v)\in\Lambda$. For such $v$, $\pi_{X\setminus Z}(v)=v-\pi_Z(v)\in\Lambda$, so $v\in(\Lambda\cap V_Z)\oplus(\Lambda\cap V_{X\setminus Z})$, and the converse inclusion is clear.

(ii) If some $\delta_x=\infty$ there is nothing to prove. Otherwise choose $a_x\in\mathfrak c_x$ with $v(a_x)=\delta_x$ and $t^{(x)}\in T$ with $t^{(x)}_x=a_x$ and $t^{(x)}_y=0$ for $y\ne x$. For $v\in\Lambda$, $t^{(x)}v=a_x\pi_x(v)$ lies in $\Lambda\cap V_x$. Hence
\[
\bigoplus_{x\in Z}a_x\pi_x(\Lambda)\ \subset\ \bigoplus_{x\in Z}(\Lambda\cap V_x)\ \subset\ M^{\rm low}\ \subset\ M^{\rm up}\ \subset\ \bigoplus_{x\in Z}\pi_x(\Lambda).
\]
Each $\pi_x(\Lambda)$ is a lattice in $V_x$, free of rank $\dim V_x$, so $\len(\pi_x(\Lambda)/a_x\pi_x(\Lambda))=\dim V_x\cdot\delta_x$. Therefore $\len(M^{\rm up}/M^{\rm low})\le\len\bigl(\bigoplus_{x\in Z}\pi_x(\Lambda)/\bigoplus_{x\in Z}a_x\pi_x(\Lambda)\bigr)=\sum_{x\in Z}\dim V_x\cdot\delta_x$.
\end{proof}

Proposition~\ref{lk:prop:ident} identifies this length with the summand in Definition~\ref{r5:def:eta} after division by $\len(O/l)$.

\begin{lemma}\label{w4:lem:det}
Let $Q\in M_n(O)$. Then $v(\det Q)=\len_O\operatorname{coker}Q$, and if $\det Q\ne0$, $\len_O\operatorname{coker}Q=\len_O\ker\bigl(Q\colon(K/O)^n\to(K/O)^n\bigr)$. More generally, if $P\colon O^m\to O^k$ is injective after $\otimes K$, then $\ker(P\otimes K/O)\cong\operatorname{coker}(P)_{\rm tors}$.
\end{lemma}

\begin{proof}
By the elementary divisor theorem $Q=U\operatorname{diag}(\varpi^{a_1},\dots,\varpi^{a_n})U'$ with $U,U'$ invertible; both sides of the first two equalities equal $\sum a_i$ (with $\infty$ allowed). For the last statement, tensor $0\to O^m\to O^k\to\operatorname{coker}P\to0$ with $K/O$: since $O^k$ is flat, $\ker(P\otimes K/O)=\operatorname{Tor}_1^O(\operatorname{coker}P,K/O)$, which is the torsion submodule of $\operatorname{coker}P$.
\end{proof}

Let $R'$ be an $O$-algebra, not necessarily commutative, and $P\in M_{a\times b}(R')$. For an $O$-algebra map $\rho\colon R'\to M_d(O)$ write $\rho(P)\in M_{ad\times bd}(O)$ for the block matrix acting $O^{bd}=(O^d)^b\to(O^d)^a$. Then $\Hom_{R'}(\operatorname{coker}(R'^a\xrightarrow{\cdot P}R'^b),V)=\ker(\rho(P)\colon V^b\to V^a)$ for $V=K^d$ or $(K/O)^d$ with $R'$ acting through $\rho$, since an $R'$-linear map from $R'^b$ is given by the images $u_1,\dots,u_b$ of the basis, and it kills the image of $\cdot P$ exactly when $\sum_jP_{ij}u_j=0$ for all $i$.

\begin{proposition}[Fitting inequality]\label{w4:prop:fitting}
Let $T\subset\prod_{x\in X}O$ and $\mathfrak c_x$, $\delta_x$ be as above. Let $R'$ be an $O$-algebra with $O$-algebra maps $\rho_x\colon R'\to M_{d_x}(O)$ for $x\in X$, and $P\in M_{a\times b}(R')$ with $a\ge b$. Assume:
\begin{enumerate}[leftmargin=*,label=(F\arabic*)]
\item for every $Q\in M_b(R')$ the function $x\mapsto\det\rho_x(Q)$ lies in $T$;
\item for every $x\ne x_0$, $\ker(\rho_x(P)\otimes K)\ne0$, that is, there is a nonzero $R'$-linear map $\operatorname{coker}P\to K^{d_x}$;
\item in a suitable $O$-basis, $\rho_{x_0}$ is block upper triangular with diagonal blocks $\rho_1,\dots,\rho_g$, $\rho_k\colon R'\to M_{n_k}(O)$, and the map $R'\to\prod_kM_{n_k}(k_O)$, $r\mapsto(\bar\rho_k(r))_k$, is surjective, where $k_O$ is the residue field;
\item $\lvert k_O\rvert>b\sum_kn_k$.
\end{enumerate}
Then
\[
\delta_{x_0}\ \le\ \sum_{k=1}^g\len_O\Hom_{R'}\bigl(\operatorname{coker}P,(K/O)^{n_k}\bigr)=\sum_{k=1}^g\len_O\ker\bigl(\rho_k(P)\ \text{on}\ ((K/O)^{n_k})^b\bigr).
\]
\end{proposition}

\begin{proof}
If some term on the right is infinite there is nothing to prove, so assume all are finite. Then each $\rho_k(P)\otimes K$ is injective: a nonzero $K$-kernel vector, scaled, gives an injection of $K/O$ into the kernel on $(K/O)^{bn_k}$, of infinite length. Let $W_k\subset O^{an_k}$ be the saturation of the image of $\rho_k(P)\colon O^{bn_k}\to O^{an_k}$, a direct summand of rank $bn_k$.

\emph{Choice of $C$.} For $C\in M_{b\times a}(R')$ put $Q=CP\in M_b(R')$. Then $\rho_k(Q)=\rho_k(C)\rho_k(P)$. There is $C$ such that for every $k$ the map $\rho_k(C)\colon O^{an_k}\to O^{bn_k}$ restricts to an isomorphism $W_k\to O^{bn_k}$. By Nakayama it suffices that $\bar\rho_k(C)$ restricts to an isomorphism $\overline W_k\to k_O^{bn_k}$, where $\overline W_k=W_k/\varpi W_k\subset k_O^{an_k}$ has dimension $bn_k$ ($W_k$ is a direct summand). By (F3), as $C$ ranges over $M_{b\times a}(R')$ the tuple $(\bar\rho_k(C))_k$ ranges over all of $\prod_kM_{bn_k\times an_k}(k_O)$, choosing each entry of $C$ independently. Fix bases of the $\overline W_k$. The function $\Phi(Y_1,\dots,Y_g)=\prod_k\det(Y_k|_{\overline W_k})$ is a polynomial on $\prod_kM_{bn_k\times an_k}(k_O)$ of total degree $\sum_kbn_k$. It is not the zero polynomial, because for each $k$ a matrix $Y_k$ that maps $\overline W_k$ isomorphically onto $k_O^{bn_k}$ exists over $k_O$. A nonzero polynomial of degree $D$ over a field with more than $D$ elements has a non-zero value at a point with coordinates in that field (induction on the number of variables). By (F4) there is such a point, and by (F3) it is $(\bar\rho_k(C))_k$ for some $C$.

\emph{The function.} Let $f(x)=\det\rho_x(Q)$. By (F1), $f\in T$. For $x\ne x_0$, $\ker\rho_x(Q)\otimes K\supseteq\ker\rho_x(P)\otimes K\ne0$ by (F2), so $f(x)=0$. Hence $f(x_0)\in\mathfrak c_{x_0}$ and $\delta_{x_0}\le v(f(x_0))$.

\emph{The value at $x_0$.} After reordering the basis, $\rho_{x_0}(Q)$ is block upper triangular with diagonal blocks $\rho_k(Q)$, so $f(x_0)=\prod_k\det\rho_k(Q)$. For each $k$, $\rho_k(C)$ maps $\operatorname{im}\rho_k(P)\subset W_k$ into $O^{bn_k}$ with $O^{bn_k}/\rho_k(C)\operatorname{im}\rho_k(P)\cong W_k/\operatorname{im}\rho_k(P)=\operatorname{coker}(\rho_k(P))_{\rm tors}$. By Lemma~\ref{w4:lem:det}, $v(\det\rho_k(Q))=\len\operatorname{coker}\rho_k(Q)=\len\operatorname{coker}(\rho_k(P))_{\rm tors}=\len\ker(\rho_k(P)\otimes K/O)$. Summing over $k$ gives the claim.
\end{proof}

In the application $R'$ is the algebra generated by the Galois group acting on an integral model of the representation attached to an eigensystem, $P$ presents the universal cocycle module of the reducible representation $\rho_\pi\oplus\bigoplus_c\chi^{c}$ with its local relations, and (F2) is a nonzero map of Ribet type at every eigensystem outside the congruence class. The application of Proposition~\ref{w4:prop:fitting} also needs a descent along a map of Hecke algebras and the joint residual surjectivity (F3) for the blocks $\rho_\pi,\chi^1,\dots,\chi^{t_r}$. \notgiven{The descent, the verification of (F2) and of (F3) for the blocks that occur, and step three are not proved here.} Step three consists of the following statements.
\begin{enumerate}[leftmargin=*,itemsep=2pt]
\item \emph{Block bounds.} For a rank one block with character $\chi^c$, the length of the space of homomorphisms from the cocycle module to $K/O$ is at most $\len\Sha_l(E,c-1)$ plus local defects at the places above $l$, at the auxiliary places and at the column places. For the block $\chi^1$ the Selmer term is $v_l(h_E)$.
\item \emph{The defect at $l$.} A lifting theorem for crystalline classes along surjections of crystalline lattices with Hodge--Tate weights in $[-h,h]$, proved with Wach modules, whose loss is $O(h/(l-1)^{j+1}+g\ell_{\max}+\log_lh)$ with main term $hl/(l-1)^2$.
\item \emph{The twist.} A ray class character of auxiliary conductor that removes the block $\rho_\pi$.
\item \emph{Twisted theta series.} Rationality of twisted theta series, their $l$-integrality with bounded denominators, the action of Hecke operators on integral models, and a reciprocity law of Shimura type.
\item \emph{The sums} over the eigensystems $\pi$, the places $w$ and the primes $l$, including the rows at $2$ and $3$.
\end{enumerate}
The published inputs of these statements include the Galois representations attached to regular algebraic polarized cuspidal automorphic representations of $\mathrm{GL}_n$ over CM fields with local--global compatibility at all places~\cite{blggt14}, the local Langlands correspondence for $\mathrm{GL}_n$, Poitou--Tate duality and the Euler characteristic formulas~\cite{nsw}, the theory of Bloch and Kato for $H^1_f$, Fontaine's period rings, Wach modules and Berger's description of Wach lattices, the Herr complex, Berger's theorem that de Rham representations are potentially semistable~\cite{berger02}, class field theory and the Chebotarev density theorem.

\subsubsection*{The local Hecke setting}

Throughout this subsection $O$ is a complete discrete valuation ring with fraction field $K$, uniformizer $\varpi$, normalized valuation $v$ and residue field $k_O$, and $l$ is the residue characteristic, so that $v(l)=\len_O(O/lO)$. We write $e_O=v(l)$. Let $X$ be a finite set, let $\mathbb T\subset\prod_{x\in X}O$ be an $O$-subalgebra which separates the points of $X$, let $M$ be a free $O$-module of finite rank with an $O$-algebra map $\mathbb T\to\operatorname{End}_O(M)$, and put $V=M\otimes_OK$, in which $M$ is a lattice. For $x\in X$ put $V_x=\{u\in V:tu=t(x)u\ \text{for all }t\in\mathbb T\}$ and $\operatorname{mult}(x)=\dim_KV_x$. For $Y\subset X$ write $V_Y=\sum_{x\in Y}V_x$. For $Y\subset Z\subset X$ put
\[
\mathfrak c_Z=\{t\in\mathbb T:t(y)=0\text{ for all }y\notin Z\},\qquad \kappa_Z(x)=\min\{v(t(x)):t\in\mathfrak c_Z\}\in\mathbb N\cup\{\infty\},
\]
\[
M^{\rm up}_Y=V_Y\cap(M+V_{X\setminus Y}),\qquad M^{\rm low}_Y=M\cap V_Y,\qquad \eta(M,Y)=\frac{\len_O(M^{\rm up}_Y/M^{\rm low}_Y)}{e_O}.
\]
The ideal $\mathfrak c_{\{x\}}$ is the congruence ideal $\mathfrak c_x$ of Proposition~\ref{w4:prop:cong}, and $\kappa_Z(x)\le\kappa_{\{x\}}(x)=\delta_x$ for $x\in Z$.

\begin{lemma}\label{lk:lem:split}
There are $C\in O\setminus\{0\}$ and $e_x\in\mathbb T$ for $x\in X$ with $e_x(y)=C$ if $y=x$ and $e_x(y)=0$ otherwise. The operators $P_x=C^{-1}e_x$ on $V$ are orthogonal idempotents with $\sum_xP_x=1$ and $P_xV=V_x$. In particular $V=\bigoplus_{x\in X}V_x$, and for $Y\subset X$ the projection $\pi_Y=\sum_{x\in Y}P_x$ of $V$ onto $V_Y$ along $V_{X\setminus Y}$ commutes with $\mathbb T$.
\end{lemma}

\begin{proof}
For $x\ne y$ choose $t\in\mathbb T$ with $t(x)\ne t(y)$ and put $s_{xy}=t-t(y)\cdot1\in\mathbb T$, so that $s_{xy}(y)=0\ne s_{xy}(x)$. The product $f_x=\prod_{y\ne x}s_{xy}\in\mathbb T$ has $f_x(x)=c_x\ne0$ and vanishes off $x$. With $C=\prod_xc_x$ the elements $e_x=(C/c_x)f_x$ lie in $\mathbb T$ and have the stated values. As functions on $X$, $e_xe_y=\delta_{xy}Ce_x$, $\sum_xe_x=C\cdot1$ and $te_x=t(x)e_x$ for $t\in\mathbb T$. Applying the algebra map to $\operatorname{End}_K(V)$ gives $P_xP_y=\delta_{xy}P_x$, $\sum_xP_x=1$, and $tP_x=t(x)P_x$, so $P_xV\subset V_x$. Conversely $P_x$ acts on $V_x$ as $C^{-1}e_x(x)=1$, so $V_x\subset P_xV$.
\end{proof}

\begin{lemma}\label{lk:lem:proj}
For $Y\subset X$, $M^{\rm up}_Y=\pi_Y(M)$. Hence $M^{\rm up}_Y/M^{\rm low}_Y$ is the module $M^{\rm up}/M^{\rm low}$ of Proposition~\ref{w4:prop:cong} for $Z=Y$ and the decomposition of Lemma~\ref{lk:lem:split}.
\end{lemma}

\begin{proof}
If $u\in V_Y$ and $u=\lambda+b$ with $\lambda\in M$ and $b\in V_{X\setminus Y}$, then $u=\pi_Y(u)=\pi_Y(\lambda)$. Conversely $\pi_Y(\lambda)=\lambda-\pi_{X\setminus Y}(\lambda)\in(M+V_{X\setminus Y})\cap V_Y$.
\end{proof}

\subsubsection*{Identification of the congruence defects}

Let $w$ be a twisted place, $l=p_w$, $1\le j\le m$, and let $\Lambda$ and $\mathbb T$ be as in Definition~\ref{r5:def:eta}, with $\mathcal O$ enlarged, if necessary, by a finite unramified extension so that every eigensystem of $\mathbb T$ on $\Lambda\otimes\mathbb Q$ has values in $\mathcal O$. This changes no length in Definition~\ref{r5:def:eta}, because for a finite unramified extension $\mathcal O'$ of $\mathcal O$ and a finite length $\mathcal O$-module $N$ one has $\len_{\mathcal O'}(N\otimes_{\mathcal O}\mathcal O')=\len_{\mathcal O}N$, and intersections and sums of submodules commute with the flat base change $\otimes_{\mathcal O}\mathcal O'$.

\begin{proposition}\label{lk:prop:ident}
The algebra $\mathbb T$ is the product $\prod_{\mathfrak m}\mathbb T_{\mathfrak m}$ of its localizations at its maximal ideals, with idempotents $\epsilon_{\mathfrak m}\in\mathbb T$, and $\Lambda=\bigoplus_{\mathfrak m}\Lambda_{\mathfrak m}$ with $\Lambda_{\mathfrak m}=\epsilon_{\mathfrak m}\Lambda$. For each $\mathfrak m$ let $X_{\mathfrak m}$ be the set of eigensystems of $\mathbb T$ on $\epsilon_{\mathfrak m}(\Lambda\otimes\mathbb Q)$, so that $\mathbb T_{\mathfrak m}$ is an $\mathcal O$-subalgebra of $\prod_{x\in X_{\mathfrak m}}\mathcal O$ separating points, and let $G_{\mathfrak m}\subset X_{\mathfrak m}$ be the eigensystems which occur in the lifts of level $j$. Then
\[
\Lambda_{\mathfrak m}=\epsilon_{\mathfrak m}(\Lambda\otimes\mathbb Q)\cap\big(\Lambda+(1-\epsilon_{\mathfrak m})(\Lambda\otimes\mathbb Q)\big),\qquad
\eta_j(w)=\sum_{\mathfrak m}\eta(\Lambda_{\mathfrak m},G_{\mathfrak m}),
\]
and each summand is the normalized length of the module $M^{\rm up}/M^{\rm low}$ of Proposition~\ref{w4:prop:cong} for $(\mathcal O,X_{\mathfrak m},\mathbb T_{\mathfrak m},\Lambda_{\mathfrak m})$ and $Z=G_{\mathfrak m}$.
\end{proposition}

\begin{proof}
The algebra $\mathbb T$ is finite over the complete local ring $\mathcal O$, so it is the product of its localizations at its finitely many maximal ideals, and the corresponding idempotents $\epsilon_{\mathfrak m}$ lie in $\mathbb T$ (idempotents lift along $\mathbb T\to\mathbb T/\varpi\mathbb T$ by completeness). Since $\epsilon_{\mathfrak m}\in\mathbb T$ preserves $\Lambda$, $\Lambda=\bigoplus_{\mathfrak m}\epsilon_{\mathfrak m}\Lambda$. If $u=\epsilon_{\mathfrak m}u=\lambda+(1-\epsilon_{\mathfrak m})u'$ with $\lambda\in\Lambda$, applying $\epsilon_{\mathfrak m}$ gives $u=\epsilon_{\mathfrak m}\lambda\in\Lambda_{\mathfrak m}$, and conversely $\epsilon_{\mathfrak m}\lambda=\lambda-(1-\epsilon_{\mathfrak m})\lambda$. This is the first formula. In Definition~\ref{r5:def:eta} the summand at $\mathfrak m$ is the normalized length of $(V_{\rm good}\cap(\Lambda_{\mathfrak m}+V_{\rm bad}))/(\Lambda_{\mathfrak m}\cap V_{\rm good})$ with $V_{\rm good}=V_{G_{\mathfrak m}}$ and $V_{\rm bad}=V_{X_{\mathfrak m}\setminus G_{\mathfrak m}}$ inside $\Lambda_{\mathfrak m}\otimes\mathbb Q$, which is $\eta(\Lambda_{\mathfrak m},G_{\mathfrak m})$ by definition, with $e_{\mathcal O}=\len_{\mathcal O}(\mathcal O/l)$. Lemma~\ref{lk:lem:proj} identifies the module with $M^{\rm up}/M^{\rm low}$.
\end{proof}

\begin{proposition}[the congruence bound at one maximal ideal]\label{lk:prop:hlb}
In the local Hecke setting let $Z'\subset Z\subset X$, $A\in\mathbb Z\setminus\{0\}$, and let $\alpha\in\operatorname{End}_O(M)$ act on $V_x$ as multiplication by $A$ for $x\in Z'$ and as $0$ for $x\in Z\setminus Z'$. Then
\[
\eta(M,Z')\ \le\ \sum_{x\in Z'}\operatorname{mult}(x)\Big(\frac{\kappa_Z(x)}{e_O}+v_l(A)\Big).
\]
\end{proposition}

\begin{proof}
If $\kappa_Z(x)=\infty$ for some $x\in Z'$ there is nothing to prove. Extending scalars to a finite unramified extension of $O$ changes neither side, as explained before Proposition~\ref{lk:prop:ident}: lengths are preserved, $M^{\rm up}_{Z'}$ and $M^{\rm low}_{Z'}$ commute with the flat base change, $\mathfrak c_Z$ is the kernel of the restriction $\mathbb T\to\prod_{y\notin Z}O$ and so commutes with it too, and the minimum of the valuations of $O'$-linear combinations of the values $t(x)$, $t\in\mathfrak c_Z$, is $\kappa_Z(x)$. So we may assume $\lvert k_O\rvert>\lvert Z'\rvert$.

\emph{A uniform element.} For $x\in Z'$ choose $t_x\in\mathfrak c_Z$ with $v(t_x(x))=\kappa_Z(x)$, and write $t_y(x)=\varpi^{\kappa_Z(x)}u_{yx}$ with $u_{yx}\in O$, so that $u_{xx}\in O^\times$. For $\lambda\in O^{Z'}$ the element $t=\sum_y\lambda_yt_y\in\mathfrak c_Z$ has $t(x)=\varpi^{\kappa_Z(x)}\sum_y\lambda_yu_{yx}$. For each $x$ the reduction of $\lambda\mapsto\sum_y\lambda_yu_{yx}$ is a nonzero linear form on $k_O^{Z'}$, and a vector space over $k_O$ is not the union of $\lvert Z'\rvert<\lvert k_O\rvert$ proper hyperplanes. So there is $\lambda$ with $v(t(x))=\kappa_Z(x)$ for every $x\in Z'$.

\emph{The operator.} Put $\beta=\alpha\circ t$ on $V$. For $\lambda\in M$, $t\lambda\in M$ and $t\lambda=\sum_{x\in Z}t(x)P_x\lambda\in V_Z$ by Lemma~\ref{lk:lem:split}, so $\alpha t\lambda=A\sum_{x\in Z'}t(x)P_x\lambda$ lies in $M\cap V_{Z'}$. The same formula gives $\beta(\pi_{Z'}\lambda)=A\sum_{x\in Z'}t(x)P_x\lambda$. With Lemma~\ref{lk:lem:proj},
\[
\beta\,M^{\rm up}_{Z'}\subset M^{\rm low}_{Z'}\subset M^{\rm up}_{Z'} .
\]
The module $M^{\rm up}_{Z'}=\pi_{Z'}(M)$ is a lattice in $V_{Z'}$, and $\beta$ acts on $V_{Z'}$ diagonally, as $At(x)$ on $V_x$, with determinant $\prod_{x\in Z'}(At(x))^{\operatorname{mult}(x)}\ne0$. Lemma~\ref{w4:lem:det} gives
\[
\len_O\big(M^{\rm up}_{Z'}/M^{\rm low}_{Z'}\big)\le\len_O\big(M^{\rm up}_{Z'}/\beta M^{\rm up}_{Z'}\big)=\sum_{x\in Z'}\operatorname{mult}(x)\big(v(A)+\kappa_Z(x)\big).
\]
Divide by $e_O$ and use $v(A)=e_O\,v_l(A)$ for $A\in\mathbb Z$.
\end{proof}

With $\alpha=1$, $A=1$ and $Z'=Z$ Proposition~\ref{lk:prop:hlb} gives $\eta(M,Z)\le e_O^{-1}\sum_{x\in Z}\operatorname{mult}(x)\kappa_Z(x)$, which refines Proposition~\ref{w4:prop:cong}(ii) because $\kappa_Z(x)\le\delta_x$.

\begin{proposition}[removing eigensystems]\label{lk:prop:cut}
In the local Hecke setting let $G\subset S\subset X$, $A\in\mathbb Z\setminus\{0\}$, and let $\alpha\in\operatorname{End}_O(M)$ act on $V_x$ as $A$ for $x\in G$ and as $0$ for $x\notin S$. Let $\bar M=M/(M\cap V_{X\setminus S})$, identified by $\pi_S$ with the lattice $\pi_S(M)$ of $V_S$, with the action of the restriction of $\mathbb T$ to $S$. Then $M\cap V_{X\setminus S}=\{m\in M:tm=0\text{ for all }t\in\mathfrak c_S\}$, and
\[
\eta(M,G)\le\eta(\bar M,G)+v_l(A)\sum_{x\in G}\operatorname{mult}(x).
\]
\end{proposition}

\begin{proof}
An element $t\in\mathfrak c_S$ kills $V_{X\setminus S}$, and $t_S=\sum_{x\in S}e_x\in\mathfrak c_S$ acts as $C\pi_S$, so $m\in M$ is killed by $\mathfrak c_S$ exactly when $\pi_Sm=0$. Since $\pi_G=\pi_G\pi_S$, $M^{\rm up}_G=\pi_G(M)=\pi_G(\bar M)=\bar M^{\rm up}_G$, and $M\cap V_G\subset\bar M\cap V_G$. If $u\in\bar M\cap V_G$, say $u=\pi_S\lambda$ with $\lambda\in M$, then $\lambda=u+\pi_{X\setminus S}\lambda$ and $\alpha\lambda=Au\in M\cap V_G$. So $A(\bar M\cap V_G)\subset M\cap V_G\subset\bar M\cap V_G$, and since $\bar M\cap V_G$ is a lattice in $V_G$,
\[
\len_O\frac{M^{\rm up}_G}{M\cap V_G}=\len_O\frac{\bar M^{\rm up}_G}{\bar M\cap V_G}+\len_O\frac{\bar M\cap V_G}{M\cap V_G}\le\len_O\frac{\bar M^{\rm up}_G}{\bar M\cap V_G}+v(A)\dim_KV_G .
\]
Divide by $e_O$.
\end{proof}

\begin{lemma}[descent along a map of Hecke algebras]\label{lk:lem:descent}
Let $\mathbb T\subset\prod_{x\in X}O$ and $\mathbb T'\subset\prod_{x'\in X'}O$ be $O$-subalgebras, $p\colon X'\to X$ a map, $Z\subset X$, and $b_1,\dots,b_q\colon X'\to O$. Assume that every $f'\in\mathbb T'$ can be written $f'=\sum_k(a_k\circ p)\,b_k$ with $a_k\in\mathbb T$, and that for $x\notin Z$ the restrictions of $b_1,\dots,b_q$ to $p^{-1}(x)$ are linearly independent over $O$. If $f'\in\mathbb T'$ vanishes on $p^{-1}(X\setminus Z)$, then for every $\pi'\in X'$, $f'(\pi')$ lies in the ideal of $O$ generated by $\{c(p(\pi')):c\in\mathfrak c_Z\}$, and $v(f'(\pi'))\ge\kappa_Z(p(\pi'))$.
\end{lemma}

\begin{proof}
Write $f'=\sum_k(a_k\circ p)b_k$. For $x\notin Z$ and $x'\in p^{-1}(x)$, $\sum_ka_k(x)b_k(x')=f'(x')=0$, so $a_k(x)=0$ for all $k$ by the independence. Hence $a_k\in\mathfrak c_Z$ and $f'(\pi')=\sum_ka_k(p(\pi'))b_k(\pi')$. For the valuation, $v(\sum_ir_iy_i)\ge\min_iv(y_i)$ in a discrete valuation ring.
\end{proof}

\begin{proposition}[from the trimmed data to Hypothesis~\ref{w4:thm:linkn}]\label{lk:prop:gd}
Fix an admissible datum and $1\le j\le m$. For each twisted place $w$ let the local components $i$ at $(j,w)$ be the maximal ideals of Proposition~\ref{lk:prop:ident}, with good sets $G_i$, and suppose that at each $i$ with $G_i\ne\emptyset$ there are $S_i\supset G_i$ and $\alpha_i$ as in Proposition~\ref{lk:prop:cut} with $A=\lvert\Gamma\rvert$ an integer independent of $i$ and $w$. Write $\eta^{\rm cut}_j(w)=\sum_i\eta(\bar\Lambda_i,G_i)$. Assume
\begin{enumerate}[leftmargin=*,label=(\alph*)]
\item $\sum_i\sum_{x\in G_i}\operatorname{mult}(x)\le d^0_j$ for every twisted $w$;
\item $\log\lvert\Gamma\rvert\le C_1x^{C_1}L_\varepsilon$, $x^{C_1}\le b$, $L_\varepsilon\ge1$ and $f,m\ge1$;
\item $\sum_{w\ \mathrm{twisted}}\ell_we_w\,\eta^{\rm cut}_j(w)\le d^0_j\big(C(fmbL_\varepsilon^2+(\log x)^2N_c)+\sum_l\log l\sum_{1\le e<t_r}\len\Sha_l(E,e)\big)$.
\end{enumerate}
Then the inequality of Hypothesis~\ref{w4:thm:linkn} holds with the constant $C+C_1$.
\end{proposition}

\begin{proof}
At a component with $G_i=\emptyset$ both $\eta(\Lambda_i,G_i)$ and $\eta(\bar\Lambda_i,G_i)$ vanish. Propositions~\ref{lk:prop:ident} and~\ref{lk:prop:cut} and (a) give $\eta_j(w)\le\eta^{\rm cut}_j(w)+d^0_j\,v_{p_w}(\lvert\Gamma\rvert)$. With $\ell_w=\log N\mathfrak p_w/[\bK:\mathbb Q]=f_w\log p_w/[\bK:\mathbb Q]$ and $\sum_{w\mid p}e_wf_w=[\bK:\mathbb Q]$,
\[
\begin{aligned}
\sum_{w\ \mathrm{twisted}}\ell_we_w\,v_{p_w}(\lvert\Gamma\rvert)&\le\sum_pv_p(\lvert\Gamma\rvert)\log p=\log\lvert\Gamma\rvert\\ &\le C_1x^{C_1}L_\varepsilon\le C_1bL_\varepsilon^2\le C_1\big(fmbL_\varepsilon^2+(\log x)^2N_c\big),
\end{aligned}
\]
because $N_c\ge0$. Multiply by $d^0_j$ and add (c).
\end{proof}

\subsubsection*{The rank $d^0_j$}

Let $\theta_j\colon\widetilde M_{\mu_0,j}\to M_j$ be the theta map from the forms of level $j$ and standard type to the space $M_j$ of level $j$, let $\theta_j^*$ be its adjoint, and let $d^0_j=\dim Z_j$ with $Z_j=\theta_j^*(M_j)$, the space of lifted forms of level $j$ and standard type. Let $D_0$ be the number of base rows, so that $D_0=\sum_\sigma n^0_\sigma$ over the isotypic types $\sigma$ with multiplicities $n^0_\sigma$.

\begin{proposition}\label{lk:prop:d0}
Suppose that $\dim M_j=D_0$ for $1\le j\le m$. Then $d^0_j\le D_0$. If moreover $\theta_j$ is surjective, then $d^0_j=D_0$. Surjectivity holds when the level projection at the standard type is surjective, $\theta_j$ restricted to its image is the composite of that projection, a bijection $\theta_{j,t}$ of coordinate spaces and a coordinate isomorphism onto $M_j$.
\end{proposition}

\begin{proof}
$d^0_j=\dim\theta_j^*(M_j)\le\dim M_j=D_0$. If $\theta_j$ is surjective, $\theta_j^*$ is injective, because $\langle\theta_j^*z,u\rangle=\langle z,\theta_ju\rangle$ vanishes for all $u$ only if $z=0$, and then $d^0_j=\dim M_j$. The last statement is the composition of surjective maps.
\end{proof}

Theorem~\ref{r5:thm:DIVinN} only needs $d^0_j\le D_0$: the bracket in Hypothesis~\ref{w4:thm:linkn} is nonnegative once $C\ge0$, which may be assumed, so the bound is monotone in $d^0_j$.

\subsubsection*{The local estimates}

Propositions~\ref{lk:prop:ident} to~\ref{lk:prop:gd} reduce Hypothesis~\ref{w4:thm:linkn} to the following statement, with $\kappa=\kappa_{Z_i}$ for the set $Z_i\supset G_i$ of eigensystems which admit a nonzero map of Ribet type, so that after the trimming of Proposition~\ref{lk:prop:cut} one has $Z_i\cap S_i=G_i$ and Proposition~\ref{lk:prop:hlb} applies with $Z'=Z=G_i$ and $\alpha=1$.

\begin{hypothesis}[per-component congruence bound]\label{lk:hyp:block}
There is $C$ such that at every component $i$ of residue characteristic $l$ at a twisted place $w$ and every $\pi\in G_i$,
\[
\frac{\kappa(\pi)}{e_O}\le s_\pi+b_l+\sum_{0<e<t_r}\len\Sha_l(E,e)+v_l(h_E)+\varepsilon_l ,
\]
where $s_\pi$ is the local term of $\rho_\pi$ and $\varepsilon_l=l^{-2}$. Away from the two auxiliary primes one also needs
\[
 s_\pi\le\nu_l+\varepsilon_l,\qquad
 \frac{t_r}{e_O}\sum_{v\in Q_i}\len_O\bigl(O/(\zeta_{i,v}-1)\bigr)\le\varepsilon_l ,
\]
where $\nu_l$ is the cyclotomic exponent of the host at $l$. At $3$ and the small primes of the host take $b_l=C(fmb\log x+x^CL_\varepsilon^2)$, and at $q,d$ take $b_l=2Cfmb$. At the other nonauxiliary primes above $l_*$ take $b_l=0$ outside an exceptional set of at most $(N_T+2)(S_c+3\eta+t_r+8b/l_*+9)$ primes $l\le3b$. For each remaining small or exceptional prime the local argument requires a length estimate $b_l\le Cfn t_r^2n_VL_\varepsilon^2+\mathcal L(l)$, where $\mathcal L(l)$ is the contribution of the split or inert local blocks. The weighted sums of these lengths over the two sets of primes also require bounds. At the two auxiliary primes $l=g_k$ one requires the bound $\eta(\bar\Lambda_i,G_i)\le\sum_{\pi\in G_i}\operatorname{mult}(\pi)r_\pi$ with
\[
\begin{gathered}r_\pi=4fn_Ve_T+2fr^2+s_\pi+\sum_{0<e<t_r}\len\Sha_l(E,e)+v_l(h_E)+4n_V\nu_i\\ {}+t_r\sum_{q\in Q_i}\frac{v(\zeta_{i,q}-1)}{e_O},\qquad e_T=C n_V^2(n_V+\eta+t_r),\end{gathered}
\]
where $\nu_i=\sum_{0<k<t_r}\big(s_{0,i}+v_l(k)\big)$ with $s_{0,i}$ the local level exponent at $l$, $n_V=r+t_r$, and $\zeta_{i,q}$ is the value of the character $\chi^1$ on a tame generator at $q\in Q_i$. Together with these, the summation of the resulting bounds over $\pi$, $i$, $w$ and $l$ gives condition (c) of Proposition~\ref{lk:prop:gd}.
\end{hypothesis}

Lemma~\ref{lk:lem:descent} and Proposition~\ref{w4:prop:fitting} bound $\kappa(\pi)$ by the Fitting lengths of the cocycle module, once (F2) and (F3) hold for the blocks $\rho_\pi,\chi^1,\dots,\chi^{t_r}$. The bounds of the Fitting lengths are the block bounds, the lifting theorem with Wach modules at the places above $l$, the twist and the statements on twisted theta series listed after Proposition~\ref{w4:prop:fitting}. \notgiven{Hypothesis~\ref{lk:hyp:block}, the verification of (F2) and (F3), and the summation are not proved here.}

For the local lifting at $l=2$, first pass to
$K_N=\mathbb Q_l(\mu_{\operatorname{lcm}(l,N)})$ with $16\mid N$.
The cyclotomic group over $K_N$ is $1+2^{v_2(N)}\mathbb Z_2$, hence is
torsion free and procyclic. Herr's cyclotomic complex applies over this
field. For $K\subset K_N$ restriction followed by corestriction is
multiplication by $[K_N:K]$, with loss $v_l([K_N:K])$ in the lifting
bound. The quantitative estimate for the crystalline lattice over $K_N$
is still required.

\subsection{The bound for the congruence defects}\label{tw:sec:divinn}

\begin{theorem}\label{r5:thm:DIVinN}
Let $0<\varepsilon<\frac1{52}$, and suppose that the admissible data at the scale $x\ge x_0$ satisfy condition (D6) of Definition~\ref{r6:R6:def:adm}, that the conclusion of Hypothesis~\ref{w4:thm:linkn} holds at $x$, that $(\log x)^2\le x^\varepsilon$, and that $d^0_j\le D_0$ for $1\le j\le m$. Then for every $1\le j\le m$
\[
\sum_{w\ \mathrm{twisted}}\ell_w\,e_w\,\eta_j(w)\le(C_N+1)\,D_0\,x^{\varepsilon}\,(B_0+N_c).
\]
\end{theorem}

\begin{proof}
Put $y=x^\varepsilon=L_\varepsilon\ge1$. Since $B_0=fmby$, the first term in the bracket of Hypothesis~\ref{w4:thm:linkn} is $fmby^2=yB_0$, and $(\log x)^2N_c\le yN_c$ by hypothesis. By (D6), $V\le\sqrt b\le b\le B_0\le y(B_0+N_c)$, where $b\le B_0$ because $fmy\ge1$. Hence the bracket is at most $C_Ny(B_0+N_c)+y(B_0+N_c)$.
\end{proof}

Theorem~\ref{r5:thm:DIVinN} enters the bound for the heights $h_u$ of the period maps $B_u$ in \eqref{A4} through the local terms $E_b$ of Hypothesis~\ref{r4:R4:HT}, and the bound for the spread of the slopes of the family $(N_u)$ in \eqref{A3}. The local terms $E_u$ of Definition~\ref{r5:def:Eterms} do not depend on $\eta_j(w)$. Proposition~\ref{lk:prop:d0} gives the required inequality once $\dim M_j=D_0$.
